% Consolidated 03-mechanics-conservation from b597304e78a782cd65c7976ce94a1a744dfe8067; original TeX fragments preserved.

% === BEGIN 03-mechanics-conservation-118-120.tex blob 6179ca54abd2a587118ffd6f9231a9151ab35983 ===
% Irodov, Arihant PDF p. 30, printed p. 21.
% Start of section 1.3: work, kinetic energy, momentum.
\exproblem{118}{21}
质点沿 $xy$ 平面中的任意轨迹从
$\bm r_1=\bm i+2\bm j$ 移动到
$\bm r_2=2\bm i-3\bm j$，
其间受一恒力 $\bm F=3\bm i+4\bm j$。
上述各量均采用 SI 单位。求该恒力做的功。
\exsolution
恒力的功不依赖所走轨迹，仅由始末位移决定：
\[
\Delta\bm r=\bm r_2-\bm r_1
=(2-1)\bm i+(-3-2)\bm j
=\bm i-5\bm j.
\]
故
\[
\boxed{W=\int_1^2\bm F\cdot\mathrm d\bm r
=\bm F\cdot(\bm r_2-\bm r_1)
=3\times1+4\times(-5)=-17\,\mathrm J}.
\]
\excheck 负功表示该外力沿净位移方向的投影为负，
并不意味着质点不能按所给轨迹运动；题目还允许其他力同时作用。

\exproblem{119}{21}
一辆质量为 $m$ 的机车从静止开始运动，
其速率满足 $v=a\sqrt s$，
$s$ 为从起点起已经走过的路程，$a$ 为常数。
求开始运动后前 $t$ 秒内全部作用力做的总功。
\exsolution
注意本版 PDF 的文本提取可能漏掉 $\sqrt s$ 的根号；
原书第21印刷页的公式经图像核验确为 $v=a\sqrt s$。
沿单向运动有
\[
\frac{\mathrm ds}{\mathrm dt}=a\sqrt s.
\]
取“自 $t=0$ 即开始运动”的非静止解，
分离变量得
$2\sqrt s=at$，于是
\[
s(t)=\frac{a^2t^2}{4},\qquad
v(t)=a\sqrt{s(t)}=\frac{a^2t}{2}.
\]
由动能定理，所有力做功之和为动能增量：
\[
\boxed{W_{\rm total}(0,t)=\frac12mv(t)^2
=\frac{ma^4t^2}{8}}.
\]
\excheck $v=a\sqrt{s}$ 与 $v(0)=0$ 单独作为微分方程
仍有“先静止、再启动”的非唯一数学解；
物理上题目“开始运动后前 $t$ 秒”已选取立即起动的分支。

\exproblem{120}{21}
质量为 $m$ 的质点沿半径为 $R$ 的圆周运动，
其动能与走过的路程 $s$ 的关系为
$K(s)=as^2$，其中 $a>0$。
求合外力的大小作为 $s$ 的函数。
\exsolution
沿运动切线，合力切向分量由功—能关系确定：
\[
F_\tau=\frac{\mathrm dK}{\mathrm ds}=2as.
\]
又 $K=\frac12mv^2$，得 $mv^2=2as^2$，
故沿圆周法线方向的合力为
\[
F_n=\frac{mv^2}{R}=\frac{2as^2}{R}.
\]
两个分量互相垂直，因而合力大小
\[
\boxed{F(s)=\sqrt{F_\tau^2+F_n^2}
=2as\sqrt{1+\frac{s^2}{R^2}}}.
\]
\excheck 不能仅把 $F_\tau=\mathrm dK/\mathrm ds$ 当作合力；
曲线运动存在向心力分量。$a$ 的量纲为
$\mathrm{J/m^2}$，所以 $2as$ 是力。
% === END 03-mechanics-conservation-118-120.tex ===

% === BEGIN 03-mechanics-conservation-121-130.tex blob e4f09d6d5a5ca4a9b43712aad51ed676c6fc6c21 ===
% Irodov, Arihant PDF pp. 30--31, printed pp. 21--22.
% The source's problem 130 has an apparent sign inconsistency;
% report both the printed relation and the physics-consistent form.
\exproblem{121}{21}
一物体质量为 $m$，被一沿路径切线方向的外力
缓慢拉上小山，始末高差为 $h$，山坡水平投影长为 $\ell$
（图1.29）。物体与山坡之间摩擦系数为 $k$。
求外力所做的功。
\exsolution
缓慢移动意味着忽略动能变化和沿切向的加速度。
在局部坡角 $\theta(s)$ 处，
支持力 $N=mg\cos\theta$，
外力切向大小为
$F=mg\sin\theta+kmg\cos\theta$。
沿轨迹积分：
\[
W=\int F\,\mathrm ds
=mg\int\sin\theta\,\mathrm ds
+kmg\int\cos\theta\,\mathrm ds.
\]
由于 $\mathrm dy=\sin\theta\,\mathrm ds$、
$\mathrm dx=\cos\theta\,\mathrm ds$，
两项积分分别为高差 $h$ 与水平投影 $\ell$，所以
\[
\boxed{W=mg(h+k\ell)}.
\]
\excheck 与山坡的具体曲线形状无关，
前提是坡面始终单调向上、摩擦阻力始终与移动反向。

\exproblem{122}{21}
质量为 $m=50\,\mathrm g$ 的圆片从静止开始
沿倾角 $\alpha=30^\circ$ 的粗糙斜面滑下，
到达水平粗糙面后再滑行 $\ell=50\,\mathrm{cm}$ 停止。
两段面的摩擦系数均为 $k=0.15$。
求整段运动中摩擦力所做的总功。
\exsolution
设斜面段走过距离为 $L$。
滑至坡底时由功—能定理
\[
K_{\rm bottom}=mgL\sin\alpha-kmgL\cos\alpha.
\]
在水平段，摩擦力做功 $-kmg\ell$，最终停止，所以
\[
mgL(\sin\alpha-k\cos\alpha)=kmg\ell
\quad\Longrightarrow\quad
L=\frac{k\ell}{\sin\alpha-k\cos\alpha}.
\]
总摩擦功为
\[
W_f=-kmg(L\cos\alpha+\ell)
=-\frac{kmg\ell}{1-k\cot\alpha}.
\]
代入数值得
\[
\boxed{W_f\simeq-4.96\times10^{-2}\,\mathrm J
\approx-0.050\,\mathrm J}.
\]
\excheck 要能顺利下滑，必须
$\tan\alpha>k$；本题满足。
负号表明摩擦耗散动能，而非创造能量。

\exproblem{123}{21}
两物块质量分别为 $m_1,m_2$，
由一根初始未形变的轻弹簧连接，置于水平粗糙面，
摩擦系数同为 $k$。
对 $m_1$ 水平施加一个恒力 $F$，
求使另一物块 $m_2$ 开始移动所需的最小 $F$。
\exsolution
设弹簧劲度系数为 $\kappa$。
在 $m_2$ 尚未运动而 $m_1$ 已向前滑动的阶段，
作用于 $m_1$ 的外力净做功为
$(F-km_1g)x$，能量分配给 $m_1$ 的动能
和弹簧弹性势能：
\[
(F-km_1g)x=\frac12m_1v_1^2+\frac12\kappa x^2.
\]
对于给定 $F$，最大弹簧形变出现在 $v_1=0$ 的转折点，
故 $\kappa x_{\max}=2(F-km_1g)$。
第二物块刚将克服静摩擦时，
弹簧拉力需达 $km_2g$。临界条件
$2(F-km_1g)=km_2g$，得到
\[
\boxed{F_{\min}=kg\left(m_1+\frac{m_2}{2}\right)}.
\]
\excheck 本题采用两块的动、静摩擦均以同一系数 $k$
描述的理想模型，且力从静止起恒定施加；
若缓慢准静态增加 $F$，阈值可能不同。

\exproblem{124}{21}
总质量 $m=0.80\,\mathrm{kg}$、
长度 $\ell=1.50\,\mathrm m$ 的均匀链条搭在粗糙桌面上，
一端悬垂。悬垂部分长度达到总长的
$\eta=1/3$ 时，链条刚可自行滑下。
求从此时起直到完全滑离桌面，
桌面对链条的摩擦力所做总功。
\exsolution
设桌面与链条之间摩擦系数为 $k$。
在临界滑动时悬垂部分重力与桌面最大摩擦力平衡：
\[
\eta mg=k(1-\eta)mg
\quad\Longrightarrow\quad
k=\frac{\eta}{1-\eta}.
\]
当悬垂长度为 $x$ 时，桌面上剩余链长为 $\ell-x$，
摩擦力大小为
$F_f(x)=k(m/\ell)(\ell-x)g$。
链条增量滑移 $\mathrm dx$ 时，桌上各微元移动
同样距离，因此摩擦力总功为
\[
W_f=-\int_{\eta\ell}^{\ell}
 k\frac m\ell(\ell-x)g\,\mathrm dx
=-\frac12kmg\ell(1-\eta)^2
=-\frac12\eta(1-\eta)mg\ell .
\]
代入得
\[
\boxed{W_f\simeq-1.31\,\mathrm J}.
\]
\excheck 积分终点时桌面上的链长为零，
摩擦力也降为零；不可假定摩擦力在整个过程中恒定。

\exproblem{125}{21}
质量为 $m$ 的物体以初速 $v_0$、
与水平夹角 $\alpha$ 斜抛，忽略空气阻力。
求：（a）重力在全飞行时间内的平均功率；
（b）重力瞬时功率随时间的变化。
\exsolution
取竖直向上为正，$v_y(t)=v_0\sin\alpha-gt$。
重力为 $\bm F_g=-mg\bm j$，故功率
\[
\boxed{P_g(t)=\bm F_g\cdot\bm v(t)
=-mgv_y(t)=mg(gt-v_0\sin\alpha)}.
\]
如果落回发射高度，飞行全程重力的总功
$W_g=-mg(y_{\rm final}-y_{\rm initial})=0$，
因此
\[
\boxed{\overline{P}_g=\frac{W_g}{T}=0}.
\]
\excheck 上升阶段 $P_g<0$，下降阶段 $P_g>0$；
全程平均值零并非瞬时功率始终为零。

\exproblem{126}{21--22}
质量为 $m$ 的质点沿半径 $R$ 的圆周运动，
其法向加速度随时间变化为 $a_n=at^2$，$a>0$。
求合外力的瞬时功率及其在最初 $t$ 秒的平均值。
\exsolution
法向加速度等于 $v^2/R$，因此在 $t\ge0$：
\[
v(t)=\sqrt{aR}\,t.
\]
运动的动能为
$K=\frac12mv^2=\frac12maRt^2$。
所有力的合功率等于动能变化率，而法向力与速度垂直不做功：
\[
\boxed{P(t)=\frac{\mathrm dK}{\mathrm dt}=maRt}.
\]
区间 $[0,t]$ 内的平均功率
\[
\boxed{\overline P(0,t)=\frac{K(t)-K(0)}t
=\frac12maRt}.
\]
\excheck 两个结果分别有比例 $t$ 与 $t/2$，
符合匀速增长功率在区间中的平均值。

\exproblem{127}{22}
一质量为 $m$ 的物块在水平地面上获得初速 $v_0$，
随后仅受滑动摩擦。
(a) 当 $k=0.27,m=1.0\,\mathrm{kg},v_0=1.5\,\mathrm{m/s}$
时，求停止前摩擦力的平均功率；
(b) 若摩擦系数随离起点距离 $x$ 变化为 $k(x)=\alpha x$，
求运动期间摩擦力瞬时功率的最大绝对值。
\exsolution
(a) 减速度大小为 $kg$，
停下用时 $t_*=v_0/(kg)$，
摩擦力总功 $W_f=-mv_0^2/2$，所以
\[
\boxed{\overline P_f=
\frac{W_f}{t_*}=-\frac12kmgv_0
\simeq-1.98\,\mathrm W}.
\]
(b) 在 $x$ 处，摩擦力 $F_f=-m\alpha gx$。
动能定理给出
\[
\frac12mv(x)^2-\frac12mv_0^2
=-\int_0^x m\alpha gu\,\mathrm du,
\quad
v(x)=\sqrt{v_0^2-\alpha gx^2}.
\]
因此瞬时功率
$P_f(x)=-m\alpha gx\sqrt{v_0^2-\alpha gx^2}$。
为使 $|P_f|$ 最大，对
$x^2(v_0^2-\alpha gx^2)$ 求极值，
可得 $x_m^2=v_0^2/(2\alpha g)$，进而
\[
\boxed{P_{f,\rm most\,negative}
=-\frac12m v_0^2\sqrt{\alpha g}},
\quad
\boxed{|P_f|_{\max}=\frac12m v_0^2\sqrt{\alpha g}}.
\]
\excheck 初始 $x=0$ 时摩擦力为零；
终点速度为零时功率也为零，
最大耗散功率发生在中间位置。

\exproblem{128}{22}
质量为 $m=0.10\,\mathrm{kg}$ 的物体位于以恒角速度
$\omega=5.0\,\mathrm{rad/s}$ 绕固定轴旋转的参考系中。
将它沿任意路径从离轴距离 $r_1=0.30\,\mathrm m$
移至 $r_2=0.50\,\mathrm m$，
求离心惯性力做的功。
\exsolution
旋转系内离心力大小为 $F_{\rm cf}=m\omega^2r$，
沿远离转轴的径向。
由于 $\bm F_{\rm cf}=m\omega^2r\,\bm e_r$，
其做功微元为
$\mathrm dW=m\omega^2r\,\mathrm dr$，
角向运动不做功。所以
\[
\boxed{W_{\rm cf}=\int_{r_1}^{r_2}
m\omega^2r\,\mathrm dr
=\frac12m\omega^2(r_2^2-r_1^2)
=0.200\,\mathrm J}.
\]
\excheck 此力在角速度恒定时可由径向有效势
$U_{\rm cf}=-m\omega^2r^2/2$ 导出，因此与路径无关。

\exproblem{129}{22}
两根轻弹簧串联，弹性系数分别为
$\kappa_1,\kappa_2$。若缓慢将其总长度拉长
$\Delta\ell$，求外力至少须做多少功。
\exsolution
串联弹簧中的拉力相同，记为 $F$，
两者伸长分别为 $x_1=F/\kappa_1$、
$x_2=F/\kappa_2$，故
\[
\Delta\ell=F\left(\frac1{\kappa_1}+\frac1{\kappa_2}\right),
\quad
\kappa_{\rm eq}=\frac{\kappa_1\kappa_2}{\kappa_1+\kappa_2}.
\]
缓慢拉伸可使末态无多余动能，
外力所做最小功即弹簧能量增量：
\[
\boxed{W_{\min}=
\frac12\kappa_{\rm eq}(\Delta\ell)^2
=\frac{\kappa_1\kappa_2}{2(\kappa_1+\kappa_2)}
(\Delta\ell)^2}.
\]
\excheck 相对较软的弹簧伸长量更大，
且等效劲度系数小于任一单根弹簧的劲度系数。

\exproblem{130}{22}
质量为 $m$ 的物体从地面被向上提升，
外力随高度 $y$ 变化。本版题面\emph{印为}
$F(y)=2(ay-1)mg$，$a>0$。
求上升过程前半段外力的功及重力势能增加量。
\exsolution
\textbf{先作源书一致性核查。}
若把印刷表达式中的 $F(y)$ 严格理解为沿\emph{向上}
的有符号外力，则在地面 $F(0)=-2mg$，合力向下，
物体从静止根本不能自发上升；而本版书末答案
$W=3mg/(4a)$、$\Delta U=mg/(2a)$ 与此符号不符。
这说明题面需要澄清外力方向或印刷符号。
为得到与书末答案一致的上升模型，明确采用
\emph{向上拉力大小}
\[
F_{\uparrow}(y)=2(1-ay)mg
\quad(0\le y\le 1/a).
\]
它与印刷式在所讨论高度范围内符号相反，
此处并未暗中改动题面，而是作为\emph{勘误后的条件分支}。
物体若由静止开始，由净力
$F_\uparrow-mg=mg(1-2ay)$ 得动能
\[
K(y)=\int_0^y mg(1-2au)\,\mathrm du
=mg(y-ay^2).
\]
它再次归零的高度是 $H=1/a$，
前半段上升对应 $0\le y\le H/2=1/(2a)$。
故该条件分支下
\[
\boxed{W_\uparrow=\int_0^{1/(2a)}2(1-ay)mg\,\mathrm dy
=\frac{3mg}{4a}},
\]
\[
\boxed{\Delta U_g=mg\frac{H}{2}=\frac{mg}{2a}}.
\]
\excheck 若\emph{完全保留题面有符号力}
$F=2(ay-1)mg$，对同样人为指定的位移区间，
积分所算外力功会是 $-3mg/(4a)$，
但这不是从静止上升到 $H$ 的自洽物理过程；
该争议须在独立审校时回查其他版本，不得标注为已解决。
% === END 03-mechanics-conservation-121-130.tex ===

% === BEGIN 03-mechanics-conservation-131-140.tex blob 5c4ab3e4b82db5456a951c30039dc95ccb33fcfa ===
% Irodov, Arihant PDF pp. 31--33, printed pp. 22--24.
\exproblem{131}{22}
某球对称力场中的势能为
$U(r)=a/r^2-b/r$，其中 $a,b>0$。
求：（a）平衡半径及其稳定性；
（b）最大吸引力，并说明 $U(r)$、径向力 $F_r(r)$ 的图像特点。
\exsolution
径向力为
\[
F_r(r)=-\frac{\mathrm dU}{\mathrm dr}
=\frac{2a}{r^3}-\frac{b}{r^2}.
\]
平衡需 $F_r=0$，故
$\boxed{r_0=2a/b}$。
在此点
\[
U''(r_0)=\left(\frac{6a}{r^4}
-\frac{2b}{r^3}\right)_{r=r_0}
=\frac{2a}{r_0^4}>0,
\]
势能极小，平衡稳定。
当 $r>r_0$ 时 $F_r<0$，吸引力大小
$f(r)=b/r^2-2a/r^3$。
求导 $f'=-2b/r^3+6a/r^4$，
最大值出现在 $r_m=3a/b$，有
\[
\boxed{f_{\max}=\frac{b^3}{27a^2}}.
\]
\excheck $r\to0^+$ 时 $U\to+\infty,F_r\to+\infty$
（强斥力）；$r\to\infty$ 时 $U\to0^-$、$F_r\to0^-$。
$U$ 在 $r_0$ 取负极小，$F_r$ 在 $r_0$ 过零，
在 $r_m$ 取最小负值。

\exproblem{132}{22}
二维势场的势能为
$U(x,y)=\alpha x^2+\beta y^2$，
其中 $\alpha,\beta>0$ 且 $\alpha\ne\beta$。
求：（a）是否为中心力场；
（b）等势线和力大小相同的点构成什么曲线，
它们的半轴比各是多少。
\exsolution
力矢量
\[
\bm F=-\nabla U=(-2\alpha x,-2\beta y).
\]
若是以原点为中心的中心力，任一点的 $\bm F$
必须与 $\bm r=(x,y)$ 平行。
但 $\bm r\times\bm F$
的平面垂直分量为 $2(\alpha-\beta)xy$，
一般不为零，因此
$\boxed{\text{非中心力场}}$。
对于给定势能 $U_0>0$，等势线
$\alpha x^2+\beta y^2=U_0$ 是椭圆；
其 $x,y$ 轴半径分别为
$\sqrt{U_0/\alpha},\sqrt{U_0/\beta}$，故
$\boxed{a_x/a_y=\sqrt{\beta/\alpha}}$。
力大小相同 $|\bm F|=F_0$ 则
\[
\alpha^2x^2+\beta^2y^2=\frac{F_0^2}{4},
\quad\boxed{a_x/a_y=\frac{\beta}{\alpha}}.
\]
\excheck 两族曲线都是同轴椭圆，但长短轴比不同；
仅 $\alpha=\beta$ 时力才处处径向。

\exproblem{133}{22}
考察两种与时间无关的平面力场：
$\bm F_1=ay\,\bm i$ 和
$\bm F_2=ax\,\bm i+by\,\bm j$，
其中 $a,b$ 均为常数。
求两力场分别是否存在势能函数。
\exsolution
在单连通平面上，足够光滑力场存在势函数的必要充分条件是
$\partial F_y/\partial x=\partial F_x/\partial y$。
对 $\bm F_1$，
\[
\frac{\partial F_{1y}}{\partial x}=0,\qquad
\frac{\partial F_{1x}}{\partial y}=a.
\]
因此当 $a\ne0$ 时不满足条件，
$\boxed{\bm F_1\text{ 非保守}}$；
$a=0$ 时它为零场，当然保守。
对 $\bm F_2$，两偏导均为零，故保守，
相应势能可取
\[
\boxed{U_2(x,y)=-\frac12ax^2-\frac12by^2+C},
\]
因为 $-\nabla U_2=\bm F_2$。
\excheck 即使力与时间无关也未必是保守力；
判断标准是旋度而不是“定常”。

\exproblem{134}{23}
质量为 $m$ 的物块以初速 $v_0$
沿倾角 $\alpha$ 的粗糙斜面向上滑动，
摩擦系数为 $k$。
求停下以前的路程及摩擦力在此过程中所做的功。
\exsolution
沿斜面向上运动时，重力切向分量及摩擦力均指向下方，
减速度大小
$a=g(\sin\alpha+k\cos\alpha)$。
由 $0=v_0^2-2as$：
\[
\boxed{s_{\rm stop}=
\frac{v_0^2}{2g(\sin\alpha+k\cos\alpha)}}.
\]
摩擦力大小 $kmg\cos\alpha$，恒反向于位移，
所以
\[
W_f=-kmg\cos\alpha\,s_{\rm stop}
=\boxed{-\frac{mv_0^2 k}
 {2(\tan\alpha+k)}}.
\]
\excheck $k=0$ 时停下距离为 $v_0^2/(2g\sin\alpha)$
且摩擦功为零；$k>0$ 时摩擦功为负。

\exproblem{135}{23}
图1.30中，小圆片从高度 $H$ 的光滑坡顶由静止滑下，
坡道末端有一个高于地面 $h$ 的水平段，
圆片从该段末端水平飞出。
求使落地点与平台边缘的水平距离 $s$ 最大的 $h$
及相应最大射程。
\exsolution
能量守恒给出水平飞出速度
$v=\sqrt{2g(H-h)}$。
从高 $h$ 的平台水平抛出，落地用时
$t=\sqrt{2h/g}$，故
\[
s(h)=vt=2\sqrt{h(H-h)}.
\]
平方为
$s^2=4[hH-h^2]=H^2-4(h-H/2)^2$，
所以在 $0<h<H$ 中，唯一最大点为
\[
\boxed{h_{\rm opt}=\frac H2,\qquad
s_{\max}=H}.
\]
\excheck $h\to0$ 或 $h\to H$ 时射程均趋零，
验证内部极大值合理。

\exproblem{136}{23}
图1.31中，小物体从高度 $h$ 由静止出发，
沿光滑槽道下滑，槽道随后变为
半径 $R=h/2$ 的竖直半圆弧。
物体在接近顶部前离开槽道，之后作抛体运动。
求离槽后自由飞行过程中到达\emph{最高点}时的速率。
\exsolution
取半圆弧的圆心距地面高度为 $R=h/2$，
用 $\theta$ 表示圆弧上质点半径与竖直向上的夹角。
其高度 $y=R(1+\cos\theta)$。
由能量守恒
\[
v^2=2g(h-y)=2gR(1-\cos\theta).
\]
在槽内向心方向列力学方程
$N+mg\cos\theta=mv^2/R$；
当 $N=0$ 时离开槽道，所以
$g\cos\theta=2g(1-\cos\theta)$，
即 $\cos\theta=2/3$、$v^2=2gR/3=gh/3$。
离开时速度沿圆弧切线，其水平方向分量为
$v_x=v\cos\theta$；之后忽略阻力，
到抛体最高点的速率即 $|v_x|$，故
\[
\boxed{v_{\rm apex}=\frac23\sqrt{\frac{gh}{3}}}.
\]
\excheck 此“最高点”是离槽后的\emph{自由抛物线顶点}，
不是圆槽最上端；二者不能混淆。

\exproblem{137}{23}
摆长 $\ell$ 的单摆起初竖直静止。
将悬点突然以某水平方向的恒速移动，
使小球相对悬点获得初速。
求让小球能够在整个竖直圆周内绳始终绷紧所需的
悬点最小速率，以及小球运动至水平位置时的绳张力。
\exsolution
设悬点移动速率为 $u$；
悬点突然运动时小球速度尚未来得及变化，
故小球相对悬点的初速大小为 $u$，
初始位置在圆周最低点。
从最低点到最高点的重力势能增量为 $2mg\ell$，
能量守恒给出
$v_{\rm top}^2=u^2-4g\ell$。
为了在圆顶保持绳绷紧，
由向心方程 $T_{\rm top}+mg=mv_{\rm top}^2/\ell$
要求 $v_{\rm top}^2\ge g\ell$。
故
\[
\boxed{u_{\min}=\sqrt{5g\ell}}.
\]
在水平位置高度较起点上升 $\ell$，
此刻速度平方
$v_{\rm side}^2=u_{\min}^2-2g\ell=3g\ell$。
重力沿指向悬点的径向分量为零，
于是
\[
\boxed{T_{\rm side}=\frac{mv_{\rm side}^2}{\ell}=3mg}.
\]
\excheck 临界最小初速对应最高点张力恰为零，
再小便会在到达最高点之前绳子松弛。

\exproblem{138}{23}
如图1.32，从上方看，一质量很小的圆片 $A$
用长为 $\ell_0$ 的水平细绳系于静止圆柱
表面切点 $B$，圆柱半径为 $R$。
初始圆片获得垂直于细绳的速度 $v_0$，
随运动细绳不断绕圆柱缠绕。
忽略摩擦，求圆片最终撞上圆柱的时间。
\exsolution
细绳张力始终垂直于圆片的瞬时速度，不做功，
所以圆片速率保持 $v_0$。
令绳子在柱面上新缠绕的圆心角为 $\varphi$，
则自由绳长
$\ell(\varphi)=\ell_0-R\varphi$。
圆片轨迹是圆柱横截面的渐开线，其弧长微元
\[
\mathrm ds=\ell(\varphi)\,\mathrm d\varphi
=(\ell_0-R\varphi)\,\mathrm d\varphi.
\]
从 $\varphi=0$ 到绳全绕完
$\varphi_*=\ell_0/R$，总路程为
\[
s=\int_0^{\ell_0/R}
(\ell_0-R\varphi)\,\mathrm d\varphi
=\frac{\ell_0^2}{2R}.
\]
故
\[
\boxed{t=\frac{s}{v_0}=\frac{\ell_0^2}{2Rv_0}}.
\]
\excheck 绳长变短时圆片的瞬时曲率半径也逐渐变小；
不能把它的运动当作以固定半径 $\ell_0$ 的圆周运动。

\exproblem{139}{23}
图1.33中，长为 $\ell$、劲度系数为 $k$ 的轻弹性绳
上端固定在 $O$，下端带轻挡块 $B$。
质量为 $m$ 的小环 $A$ 从 $O$ 处静止开始下落，
先自由落下距离 $\ell$，随后挡块受力使弹性绳伸长。
求绳子的最大伸长量 $\Delta\ell$。
\exsolution
从初始位置到最大伸长点，
小环总下降距离为 $\ell+\Delta\ell$，
初末速率均为零。
忽略绳和挡块的质量以及耗散，
重力势能损失全部转为弹性势能：
\[
mg(\ell+\Delta\ell)=\frac12 k(\Delta\ell)^2.
\]
这是关于 $\Delta\ell$ 的一元二次方程，
取正根得
\[
\boxed{\Delta\ell=
\frac{mg}{k}\left(1+
\sqrt{1+\frac{2k\ell}{mg}}\right)}.
\]
\excheck $\ell\to0$ 时退化为从自然长位置
自由下落的竖直弹簧最大伸长 $2mg/k$；
另一负根没有物理意义。

\exproblem{140}{24}
图1.34中，质量为 $m$ 的物块 $A$
在光滑水平面上，初始由左侧绳 $PA$ 固定，
向右的绳跨过轻滑轮连接等质量重物 $B$。
另有弹簧竖直连接 $A$ 与其正上方固定点 $O$，
弹簧自然长 $\ell_0=0.50\,\mathrm m$，
劲度系数 $k=5mg/\ell_0$。
烧断左绳后，求 $A$ 刚开始离开水平面时的速度。
\exsolution
设 $A$ 向右滑过 $x$，由右绳长度守恒，
$B$ 下降同样距离 $x$，两物块速率相同为 $v$。
弹簧长度变为
$L=\sqrt{\ell_0^2+x^2}$。
弹簧向上的作用力分量为
$k(L-\ell_0)\ell_0/L$；
离面临界时地面支持力 $N=0$，故
\[
k(L-\ell_0)\frac{\ell_0}{L}=mg.
\]
代入 $k=5mg/\ell_0$ 得
$5(1-\ell_0/L)=1$，
所以 $L=5\ell_0/4$，进而
$x=3\ell_0/4$。
系统重物的势能损失为 $mgx$，
弹簧势能增加为
$\tfrac12 k(L-\ell_0)^2=5mg\ell_0/32$。
两等质量物块总动能为
$2(\tfrac12 mv^2)=mv^2$，故
\[
mv^2=mg\frac{3\ell_0}{4}
-\frac{5mg\ell_0}{32}
=\frac{19mg\ell_0}{32}.
\]
因此
\[
\boxed{v=\sqrt{\frac{19g\ell_0}{32}}
\simeq1.70\,\mathrm{m/s}}.
\]
\excheck 此刻弹簧竖直分力恰好抵消重力，
而非弹簧长度刚超过自然长时立刻离面。
% === END 03-mechanics-conservation-131-140.tex ===

% === BEGIN 03-mechanics-conservation-141-150.tex blob 4b5af215b426aa313c19b62747f95c373eee6247 ===
% Irodov Arihant PDF pp. 33--35, printed pp. 24--26.
\exproblem{141}{24}
如图1.35，质量 $m=1.0\,\mathrm{kg}$ 的小块置于一块
可水平缓慢移动的木板上，借一根初长
$\ell_0=0.40\,\mathrm m$ 的轻弹性绳连接到固定点 $O$。
绳初始竖直。小块与木板摩擦系数 $k=0.20$。
缓慢向右移动木板，直至弹性绳与竖直夹角
$\theta=30^\circ$、小块即将相对木板滑动。
求截至此时摩擦力对小块在地面系中所做的功。
\exsolution
设弹性绳劲度系数为 $\kappa$，
小块水平移动 $x$，绳长
$L=\sqrt{\ell_0^2+x^2}$，
伸长 $\Delta L=L-\ell_0$。
在每个准静态位置，木板摩擦力大小等于
弹性绳水平分力，
$f=\kappa(L-\ell_0)x/L$。
因 $\mathrm dL=(x/L)\mathrm dx$，
摩擦力累计做功
\[
W_f=\int_0^{x_*}f\,\mathrm dx
=\int_{\ell_0}^{L_*}\kappa(L-\ell_0)\,\mathrm dL
=\frac12\kappa(\Delta L_*)^2.
\]
到将滑动的临界位置，支持力
$N=mg-\kappa\Delta L_*\cos\theta$，而
$f=kN$，所以
\[
\kappa\Delta L_*(\sin\theta+k\cos\theta)=kmg.
\]
又 $\Delta L_*=\ell_0(\sec\theta-1)$，
故
\[
\boxed{W_f=
\frac{kmg\ell_0(1-\cos\theta)}
 {2\cos\theta(\sin\theta+k\cos\theta)}
\simeq0.090\,\mathrm J}.
\]
\excheck 此功为正，因为地面系中木板摩擦力推动小块向右；
绳储存的弹性能恰好是该做功。
板—块间接近静止的相对位移并不等于地面系的位移。

\exproblem{142}{24}
光滑轻水平杆 $AB$ 能绕端点 $A$ 的竖直轴旋转。
一小滑套质量为 $m$，以一根自然长 $\ell_0$、
劲度系数 $\kappa$ 的轻弹簧与点 $A$ 相连。
问为使系统缓慢加速到角速度 $\omega$，
外界至少应做多少功。
\exsolution
设最终稳定径向位置为 $r$。
在旋转参考系中径向力平衡：
\[
\kappa(r-\ell_0)=m\omega^2 r.
\]
令 $\eta=m\omega^2/\kappa$。
当 $0\le\eta<1$ 时
$r=\ell_0/(1-\eta)$，
弹簧伸长
$\Delta r=\ell_0\eta/(1-\eta)$。
缓慢加速可使径向振动动能趋零，
外力至少须提供最终旋转动能和弹簧弹性势能：
\[
W_{\min}
=\frac12m\omega^2r^2
+\frac12\kappa(\Delta r)^2.
\]
代入化简
\[
\boxed{W_{\min}=
\frac{\kappa\ell_0^2}{2}
\frac{\eta(1+\eta)}{(1-\eta)^2}},
\qquad
\eta=\frac{m\omega^2}{\kappa}<1.
\]
\excheck $\eta\to0$ 时
$W_{\min}\sim \tfrac12m\omega^2\ell_0^2$；
$\eta\to1^-$ 时位移与所需功发散，
意味着理想弹簧模型没有有限稳定平衡。

\exproblem{143}{24}
轻绳跨过固定的无摩擦轻滑轮，两端挂质量
$m_1,m_2$ 的重物。求两物体构成的系统质心加速度。
\exsolution
取竖直向下为正，两端加速度大小相等、方向相反。
阿特伍德机中 $m_1$ 的加速度为
\[
a_1=g\frac{m_1-m_2}{m_1+m_2},
\quad a_2=-a_1.
\]
系统总质量 $M=m_1+m_2$，
质心加速度按定义为
\[
a_C=\frac{m_1a_1+m_2a_2}{M}
=\frac{m_1-m_2}{m_1+m_2}a_1.
\]
故
\[
\boxed{a_C=g\frac{(m_1-m_2)^2}{(m_1+m_2)^2}},
\]
沿竖直向下方向（相等质量时为零）。
\excheck 虽然一个重物向上运动，另一重物向下运动，
质心仍可能向下加速；不可把两物体的加速度简单平均。

\exproblem{144}{25}
两个相互作用的质点构成孤立系统，质心始终静止。
图1.36给出了某时刻两质点的位置，
以及质量为 $m_1$ 的质点的轨迹。
另有 $m_2=m_1/2$。求画出另一个质点的轨迹。
\exsolution
由于质心静止，它的位置矢量 $\bm R_C$ 不随时间变化：
\[
m_1(\bm r_1-\bm R_C)+
m_2(\bm r_2-\bm R_C)=\bm0.
\]
利用 $m_2=m_1/2$，得到逐时刻的位置映射
\[
\boxed{\bm r_2(t)-\bm R_C
=-2[\bm r_1(t)-\bm R_C]}.
\]
因此，以固定质心为中心将图中质点1的整条轨迹
进行\emph{中心反演}，再将相对质心的距离扩大 $2$ 倍，
即得到质点2的轨迹。任意相应时刻两质点与质心共线、
位于质心的两侧。
\excheck 本版题面英文末句重复称“画 $m_1$ 的轨迹”，
而图中已给出 $m_1$ 的轨迹；
本稿按显然的物理意图给出 $m_2$ 的轨迹，
并将这一题面指代问题列为待校勘事项。

\exproblem{145}{25}
一闭合链条质量为 $m=0.36\,\mathrm{kg}$，
由一根细绳系在竖直旋转轴上，
以恒角速度 $\omega=35\,\mathrm{rad/s}$ 绕轴转动。
细绳与竖直方向夹角为 $\theta=45^\circ$（图1.37）。
求链条质心到旋转轴的距离 $r$ 及绳张力 $T$。
\exsolution
把整个闭合链条视为整体，内部约束力合力为零。
质心受力仅有向下重力 $mg$ 和沿绳向上、向轴的张力。
竖直平衡与水平向心运动分别满足
\[
T\cos\theta=mg,\qquad
T\sin\theta=m\omega^2r.
\]
两式相除消去 $T$：
\[
\boxed{r=\frac{g\tan\theta}{\omega^2}
\simeq0.00800\,\mathrm m=0.80\,\mathrm{cm}},
\]
再由竖直平衡得
\[
\boxed{T=\frac{mg}{\cos\theta}
\simeq4.99\,\mathrm N}.
\]
\excheck 只用整个链条的质心动力学便可求此二量，
不必知道链条实际受旋转影响形成的空间形状。

\exproblem{146}{25}
一均匀圆锥 $A$ 的质量
$m=3.2\,\mathrm{kg}$、半顶角 $\alpha=10^\circ$，
沿另一圆锥面 $B$ 匀速纯滚动，锥顶 $O$ 固定（图1.38）。
$A$ 的质心与 $O$ 等高，距 $O$ 的距离
$\ell=0.17\,\mathrm m$，其轴以角速度 $\omega$ 转动。
求：（a）$\omega=1.0\,\mathrm{rad/s}$ 时静摩擦力；
（b）当摩擦系数 $k=0.25$ 时保持纯滚动的角速度范围。
\exsolution
质心绕竖直轴作水平圆周运动，加速度大小为
$\omega^2\ell$，方向指向 $O$。
把接触处作用力分解成法向支持力 $N$ 与沿接触母线的
静摩擦力 $f$。在过 $O$ 和质心的竖直剖面内，
将重力与质心的水平向心加速度分别投影到接触母线
及其法线，得
\[
\boxed{f=m(g\sin\alpha+\omega^2\ell\cos\alpha)},\qquad
\boxed{N=m(g\cos\alpha-\omega^2\ell\sin\alpha)}.
\]
当 $\omega=1.0\,\mathrm{rad/s}$，
\[
\boxed{f\simeq
3.2[9.8\sin10^\circ+0.17\cos10^\circ]
\simeq5.98\,\mathrm N}.
\]
纯滚动要求 $f\le kN$，所以
\[
g\sin\alpha+\omega^2\ell\cos\alpha
\le k(g\cos\alpha-\omega^2\ell\sin\alpha).
\]
若 $k>\tan\alpha$，其角速度上限为
\[
\boxed{0\le\omega\le
\sqrt{\frac{g(k-\tan\alpha)}
{\ell(1+k\tan\alpha)}}\simeq2.02\,\mathrm{rad/s}}.
\]
\excheck 题设 $k=0.25>\tan10^\circ$，
故静止或缓慢滚动时有足够摩擦；
该上限同时保证 $N>0$。
书末给出的“$\omega\ll2$”应理解为不超过静摩擦临界值，
而非要求速度必须远小于某数值的渐近条件。

\exproblem{147}{25}
在惯性系 $K$ 中，两质点质量分别为 $m_1,m_2$，
沿同一 $x$ 轴以速率的有符号值 $v_1,v_2$ 运动。
求：（a）使两质点总动能最小的匀速平移参考系 $K'$
相对 $K$ 的速度 $V$；
（b）该最小总动能。
\exsolution
在以 $V$ 平移的惯性系中两物体的速度为
$v_1-V$ 和 $v_2-V$，总动能
\[
K(V)=\frac12m_1(v_1-V)^2+
\frac12m_2(v_2-V)^2.
\]
对 $V$ 微分求极小值：
\[
\frac{\mathrm dK}{\mathrm dV}
=-m_1(v_1-V)-m_2(v_2-V)=0.
\]
因此
\[
\boxed{V=\frac{m_1v_1+m_2v_2}{m_1+m_2}},
\]
即质心速度。代回得
\[
\boxed{K_{\min}=
\frac12\mu(v_1-v_2)^2},\quad
\mu=\frac{m_1m_2}{m_1+m_2}.
\]
\excheck 二阶导数
$\mathrm d^2K/\mathrm dV^2=m_1+m_2>0$，
故这确实是全局最小值。

\exproblem{148}{25}
某粒子系统总质量为 $M$，
其质心静止系相对惯性系 $K$ 以速度 $\bm V$ 匀速平移。
系统在质心系中的总机械能为 $\widetilde E$，
求在 $K$ 系中的总机械能 $E$。
\exsolution
质点 $i$ 的速度写为
$\bm v_i=\bm V+\widetilde{\bm v}_i$，
且质心系满足
$\sum_i m_i\widetilde{\bm v}_i=\bm0$。
因此
\[
\sum_i\frac12m_i v_i^2
=\frac12M V^2+
\sum_i\frac12m_i\widetilde v_i^2,
\]
因为动能展开后的交叉项为零。
系统内部势能与整体平移速度无关，
故总机械能关系为
\[
\boxed{E=\widetilde E+\frac12M V^2}.
\]
\excheck 这就是柯尼希动能定理；
若“总能量”包含狭义相对论静能，则不能直接将此
非相对论经典式作为严格关系。

\exproblem{149}{25--26}
两个小圆片质量分别为 $m_1,m_2$，
通过轻弹簧相连并置于光滑水平面上，
随后分别给予互相垂直的初速度矢量
$\bm v_1,\bm v_2$。
求系统在质心系中的总能量（取弹簧初始未形变）。
\exsolution
在初始时刻，系统的质心系动能为
\[
K_{\rm rel}
=\frac12\mu|\bm v_1-\bm v_2|^2,\qquad
\mu=\frac{m_1m_2}{m_1+m_2}.
\]
由于 $\bm v_1\cdot\bm v_2=0$，
\[
|\bm v_1-\bm v_2|^2=v_1^2+v_2^2.
\]
若弹簧初始处于自然长度，
则初始弹性势能为零，总能量守恒给出
\[
\boxed{\widetilde E=
\frac12\frac{m_1m_2}{m_1+m_2}(v_1^2+v_2^2)}.
\]
\excheck 若初始弹簧已发生形变，题目需要额外给出初始弹性能，
此结果应再加该项；不能由速度数据单独决定全部机械能。

\exproblem{150}{26}
两个小球质量为 $m_1,m_2$，通过轻弹簧连接。
$t=0$ 时其初速度分别为矢量
$\bm v_1,\bm v_2$，随后在地球近地面均匀重力场中运动，
忽略空气阻力。求系统总动量 $\bm p(t)$
和质心相对初始位置的位矢 $\Delta\bm r_C(t)$。
\exsolution
系统内部弹簧力总和为零，唯一外力合力是
$(m_1+m_2)\bm g$，
因此由总动量定理
\[
\frac{\mathrm d\bm p}{\mathrm dt}=(m_1+m_2)\bm g,
\]
积分得到
\[
\boxed{\bm p(t)=m_1\bm v_1+m_2\bm v_2
 +(m_1+m_2)\bm g\,t}.
\]
质心速度为
$\bm v_C(t)=\bm p(t)/(m_1+m_2)$，
所以
\[
\boxed{\Delta\bm r_C(t)=
\frac{m_1\bm v_1+m_2\bm v_2}{m_1+m_2}\,t
+\frac12\bm g\,t^2}.
\]
\excheck 无论弹簧是否振动、两球间距如何变化，
只要外力只含均匀重力，质心仍作标准匀加速抛体运动。
% === END 03-mechanics-conservation-141-150.tex ===

% === BEGIN 03-mechanics-conservation-151-160.tex blob 9ad315095f71303d0ef6106c0b2bcff463c5f18f ===
% Irodov, Arihant PDF pp. 35--36, printed pp. 26--27.
\exproblem{151}{26}
图1.39中，两个物块质量为 $m_1,m_2$，以劲度系数
$\kappa$ 的轻弹簧相连，置于光滑水平面上。
物块1起初靠着墙，物块2将弹簧向左压缩一小段
$x$ 后由静止释放。求物块1刚离墙后系统质心的速度。
\exsolution
在弹簧尚未恢复原长之前，弹簧对物块1的压力向左，
墙的支持力使物块1始终静止。物块2单独作弹簧振动，
弹簧恢复自然长度瞬间，物块1刚要离开墙。
在此前阶段，初始弹性能
$\tfrac12\kappa x^2$ 全部变成物块2的动能：
\[
\frac12\kappa x^2=\frac12m_2v_2^2
\quad\Longrightarrow\quad
v_2=x\sqrt{\frac{\kappa}{m_2}}.
\]
此瞬间 $v_1=0$，系统质心速度为
\[
\boxed{v_C=\frac{m_2v_2}{m_1+m_2}
=\frac{x\sqrt{\kappa m_2}}{m_1+m_2}},
\]
方向远离墙。
\excheck 不能从 $t=0$ 起应用\emph{两物块}的总水平
动量守恒，因为在离墙前墙对物块1有外力。

\exproblem{152}{26}
两个物块以自然长 $\ell_0$ 的轻弹簧连接，
弹簧劲度系数 $\kappa$，放在光滑水平面上。
从 $t=0$ 起对右侧物块 $m_2$ 施加恒定水平向右力 $F$。
求两物块距离的最大值与最小值：
（a）当两物块质量相等；
（b）任意质量 $m_1,m_2$。
\exsolution
设两块坐标 $x_1,x_2$，相对距离
$\ell=x_2-x_1$，则
\[
m_1\ddot x_1=\kappa(\ell-\ell_0),\qquad
m_2\ddot x_2=F-\kappa(\ell-\ell_0).
\]
相减得到
\[
\ddot\ell+\kappa\left(\frac1{m_1}+\frac1{m_2}\right)
(\ell-\ell_0)=\frac F{m_2}.
\]
令 $\mu=m_1m_2/(m_1+m_2)$，
$\Omega^2=\kappa/\mu$，
$\delta=Fm_1/[\kappa(m_1+m_2)]$。
用初值 $\ell(0)=\ell_0,\ \dot\ell(0)=0$ 解得
\[
\ell(t)=\ell_0+\delta(1-\cos\Omega t).
\]
因此
\[
\boxed{\ell_{\min}=\ell_0,\qquad
\ell_{\max}=\ell_0+\frac{2m_1F}{\kappa(m_1+m_2)}}.
\]
当 $m_1=m_2$ 时，
$\boxed{\ell_{\max}=\ell_0+F/\kappa}$。
\excheck 任意时间 $1-\cos\Omega t\in[0,2]$，
所以给出的最大、最小距离均可以实际达到。

\exproblem{153}{26}
图1.41中两个质量均为 $m$ 的小立方体上下排列，
以劲度系数 $\kappa$ 的竖直轻弹簧相连，
弹簧初始压缩 $\Delta\ell$，两物块又用绳固定。
烧断绳后，求：（a）下方立方体最终会离地的条件；
（b）若 $\Delta\ell=7mg/\kappa$，
系统质心能上升的最大高度。
\exsolution
定义 $q=mg/\kappa$，取弹簧形变
$x$（拉伸为正）。下方立方体未离地时，上方物块
从 $x=-\Delta\ell$ 由静止运动。
其机械能守恒：
\[
\tfrac12mv^2+\tfrac12\kappa x^2+mgx
=\tfrac12\kappa(\Delta\ell)^2-mg\Delta\ell.
\]
若下块一直贴地，上块最高时 $v=0$，
非平凡转折形变为
$x_{\max}=\Delta\ell-2q$。
下块受弹簧上拉力 $\kappa x$，
只有当 $x>q$ 才可能离地，因此
\[
\boxed{\Delta\ell>3q=\frac{3mg}{\kappa}}.
\]
在刚离地时 $x=q$，上块动能由上述能量式给出
\[
v^2=\frac{\kappa}{m}(\Delta\ell+q)(\Delta\ell-3q).
\]
此刻下块速度仍为零，因此系统质心速度为 $v/2$。
从初始到下块刚离地，质心已升高
$(\Delta\ell+q)/2$；之后系统只受重力，
质心再升高 $v^2/(8g)$，故
\[
\boxed{h_C^{\max}=
\frac{(\Delta\ell+q)^2}{8q}
=\frac{mg}{8\kappa}
\left(1+\frac{\kappa\Delta\ell}{mg}\right)^2}.
\]
代入 $\Delta\ell=7q$：
$\boxed{h_C^{\max}=8mg/\kappa}$。
\excheck 判断下块离地必须比较弹簧\emph{最大拉伸}
而非初始压缩时的弹簧力。
质心在两块离地以后受的加速度始终是 $-\bm g$。

\exproblem{154}{26}
两辆相同小车1、2在平行无摩擦轨道上相向运动，
每辆空车质量为 $M$，各载一名质量 $m$ 的人。
小车相遇时两人沿垂直轨道方向跳到对方车上。
最终车1静止，车2保持原来的运动方向、
速度为 $v$。求两车原来沿轨道的有符号速度 $v_1,v_2$。
\exsolution
跳跃沿横向进行，对运动方向的碰撞可以看作
“外来乘客与车体”之间的非弹性结合，
各乘客在跳出时保留原车的纵向速度。
故最终纵向速度为
\[
V_1=\frac{Mv_1+mv_2}{M+m}=0,\quad
V_2=\frac{Mv_2+mv_1}{M+m}=v.
\]
第一式给出 $v_1=-(m/M)v_2$，
代入第二式后得
\[
\boxed{v_1=-\frac{mv}{M-m},\qquad
v_2=\frac{Mv}{M-m}}.
\]
\excheck 要让车2继续沿初始方向运动且该速度为正，
必须有 $M>m$。两车最初相向：
$v_1$ 与 $v_2$ 确为相反符号。

\exproblem{155}{27}
两辆相同空车质量为 $M$，
沿同一直线无摩擦轨道同向以速率 $v_0$ 行驶。
后车载一名质量 $m$ 的乘客，乘客以相对后车的
向前速度 $u$ 跳到前车。
求两车最终速度。
\exsolution
设乘客起跳后后车速度为 $V_r$。
乘客离后车时，其纵向速度为 $V_r+u$。
对起跳过程的后车—乘客体系守恒纵向动量：
\[
(M+m)v_0=MV_r+m(V_r+u),
\]
故
\[
\boxed{V_r=v_0-\frac{m}{M+m}u}.
\]
乘客随后落到前车上，为完全非弹性纵向结合，
落地前乘客速度为
$v_p=V_r+u=v_0+Mu/(M+m)$。
对前车及乘客守恒纵向动量：
\[
(M+m)V_f=Mv_0+mv_p,
\]
所以
\[
\boxed{V_f=v_0+\frac{mM}{(M+m)^2}u}.
\]
\excheck 后车减速、前车加速；整体纵向动量
在两次相互作用全过程仍然守恒。

\exproblem{156}{27}
一辆质量为 $M$ 的无摩擦小车上静止站着两个人，
每人质量均为 $m$。
两人分别以相对小车的水平速度 $u$
向同一方向跳下。求小车最终速度：
（1）两人同时跳；
（2）两人先后跳；
并比较两种情况下的小车速度大小。
\exsolution
取乘客跳出方向为正。
（1）同时跳出时，小车末速 $V_{\rm sim}$，
两人末速均为 $V_{\rm sim}+u$，动量守恒：
\[
0=MV_{\rm sim}+2m(V_{\rm sim}+u),
\quad
\boxed{V_{\rm sim}=-\frac{2m}{M+2m}u}.
\]
（2）先跳者离车时，仍留在车上的人
与空车总质量 $M+m$，末速为 $V_1$。
第一跳：
$0=(M+m)V_1+m(V_1+u)$，
故 $V_1=-mu/(M+2m)$。
第二跳时，由相对于第一跳前的运动系守恒动量，
空车最终速度再减少 $mu/(M+m)$，于是
\[
\boxed{V_{\rm seq}=
-\frac{m(2M+3m)}{(M+m)(M+2m)}u}.
\]
二者速度大小之比为
\[
\boxed{\frac{|V_{\rm seq}|}{|V_{\rm sim}|}
=1+\frac{m}{2(M+m)}>1}.
\]
\excheck 先后跳下时第二人面对的可反冲剩余总质量更小，
故同样相对跳速产生更大最终反冲。

\exproblem{157}{27}
一条均匀柔软链条竖直悬挂，下端恰触到桌面。
烧断上端支承绳后，链条开始自由落下并逐渐堆积在桌上。
证明在任意时刻，\emph{仍在下落的链段}对桌面的冲击力
是\emph{已经静置在桌上链段}重力的两倍。
\exsolution
设链条线密度为 $\lambda$，
在时刻 $t$ 已经到达桌面的链长为 $x$。
由于烧断后各链元由静止开始自由落体，
达到桌面的链元速率为 $v=\sqrt{2gx}$。
单位时间内到达桌面的链条质量
$\dot m=\lambda v$。
下落链元被桌面制止，需要动量变化率
\[
F_{\rm impact}=\dot m\,v=\lambda v^2=2\lambda gx.
\]
已经堆积在桌面上的链条质量为 $\lambda x$，
其静重力为
$W_{\rm rested}=\lambda gx$，因此
\[
\boxed{F_{\rm impact}=2W_{\rm rested}}.
\]
\excheck 全部对桌面的瞬时压力实际上为
$F_{\rm impact}+W_{\rm rested}=3\lambda gx$，
但题目只把\emph{下落部分的冲击力}与静置部分重量相比。
假定链元落桌后立即静止且无弹跳。

\exproblem{158}{27}
一颗钢球质量 $m=50\,\mathrm g$，
由高度 $h=1.0\,\mathrm m$ 落向很重的水平板。
钢球每次碰撞后的反弹速率
是碰撞前速率的 $1/\eta$，
其中 $\eta=1.25$。
求经无数次弹跳，钢球传给平板的总冲量。
\exsolution
第一次碰撞前的下降速率
$v_0=\sqrt{2gh}$。
第 $n$ 次（从 $n=0$ 编号）碰撞前速率
$v_n=v_0/\eta^n$，碰后向上速率 $v_n/\eta$。
按向下为正，钢球碰撞动量变化为
$-m v_n/\eta-(+m v_n)$，
所以钢球传给平板的向下冲量大小为
$m v_n(1+1/\eta)$。
总冲量是收敛几何级数：
\[
J=\sum_{n=0}^\infty m v_0
\left(1+\frac1\eta\right)\eta^{-n}
=mv_0\frac{\eta+1}{\eta-1}.
\]
代入数值：
\[
\boxed{J=
0.050\sqrt{2(9.81)(1.0)}
\frac{2.25}{0.25}
\simeq1.99\,\mathrm{N\,s}},
\]
方向向下。
\excheck 第一次碰撞的冲量约 $0.399\,\mathrm{N\,s}$，
后续总和约 $1.60\,\mathrm{N\,s}$，符合几何级数。
原书印刷答案页确实写有 $0.2\,\mathrm{kg\,m/s}$，但将题面参数代入同页公式应得约 $2.0\,\mathrm{kg\,m/s}$；因此这是印刷数值勘误，并非 PDF 文字提取错误。

\exproblem{159}{27}
质量为 $M$ 的木筏与质量为 $m$ 的人一起静止浮在湖面，
忽略水的水平阻力。
人在筏上相对木筏移动总距离 $\ell'$，
相对速度随时间变化为 $v'(t)$，最后停住。
求：（a）木筏相对湖岸的位移；
（b）人对木筏的水平作用力。
\exsolution
取人的移动方向为正，筏岸位移记 $X$，
人岸位移为 $X+\ell'(t)$。
系统初始总水平动量为零，始终没有外部水平合力：
\[
M\dot X+m(\dot X+v')=0
\Longrightarrow
\dot X=-\frac{m}{M+m}v'.
\]
积分得木筏的最终位移
\[
\boxed{\Delta X=-\frac{m}{M+m}\ell'}.
\]
对木筏单独应用牛顿第二定律，
人施加的水平作用力为
\[
\boxed{F_{\rm man\to raft}=M\ddot X
=-\frac{mM}{M+m}\frac{\mathrm dv'}{\mathrm dt}}.
\]
\excheck 人在木筏上加速前行时筏向后受力，
人在末段减速时木筏可能受向前的力；
人停下后系统质心仍保持原水平位置。

\exproblem{160}{27}
一固定滑轮上跨轻绳，一端悬挂梯子和攀登者，
另一端悬挂质量为 $M$ 的配重。
攀登者质量为 $m$，在梯子上相对梯子向上爬
$\ell'$ 后停住，忽略滑轮摩擦与绳质量。
求这个系统质心的竖直位移。
\exsolution
为了使题目无需给梯子的质量还能得到唯一结果，
应采用原书常见的\emph{初始平衡配重}模型：
梯子与人的总质量等于配重 $M$，
即梯子质量为 $M-m$。
设梯子相对地面竖直位移为 $y$（向上为正），
则另一端配重因绳长约束位移为 $-y$；
人相对梯子额外上升 $\ell'$，
故其总位移为 $y+\ell'$。
总质量为 $2M$，于是质心位移
\[
\Delta y_C=
\frac{(M-m)y+m(y+\ell')+M(-y)}{2M}
=\boxed{\frac{m\ell'}{2M}}.
\]
更一般地，若梯子的实际质量为$M_L$
而非必然等于$M-m$，并令其地面位移为$y$，
则严格质心位移是
\[
 \boxed{\Delta y_C=
 \frac{(M_L+m-M)y+m\ell'}{M_L+m+M}}.
\]
它在$M_L+m=M$的初始平衡条件下，
才\emph{独立于}未知的梯子位移$y$，
并化为上面的$m\ell'/(2M)$。
\excheck 已用用户提供的原版PDF第27印刷页
复核英文原题：确实\emph{没有单独给出梯子质量}，
也没有明确写“梯子加上攀登者与配重等重”。
因此该题\emph{原印条件不足以唯一确定数值}，
书末结果只能作为附加平衡假设下的分支解，
不得标记为无附加条件的严格唯一结论。
% === END 03-mechanics-conservation-151-160.tex ===

% === BEGIN 03-mechanics-conservation-161-170.tex blob 706931fc88458fc684a8301d0a7419b3ecc7024e ===
% Irodov, Arihant PDF pp. 36--38, printed pp. 27--29.
% Momentum conservation, ballistic pendulum, elastic/inelastic collisions.
\exproblem{161}{27}
质量为 $M$ 的火炮从静止开始沿倾角为 $\alpha$ 的
光滑斜面自由滑下，沿斜面移动距离 $\ell$ 后射出一枚
水平方向动量大小为 $p$ 的炮弹（相对于地面）。
因反冲，火炮在发射后停止。
炮弹质量相对火炮可忽略。求发射持续时间 $\tau$。
\exsolution
设火炮沿斜面向下的速度刚开始发射时为 $u$，
由发射前机械能守恒
\[
\frac12Mu^2=Mg\ell\sin\alpha
\quad\Longrightarrow\quad
u=\sqrt{2g\ell\sin\alpha}.
\]
取沿斜面向下为正。
在发射持续时间 $\tau$ 内，
重力沿斜面的冲量为 $Mg\sin\alpha\,\tau$。
为使火炮停下，炮弹水平方向发射须在斜面向下方向
具有正投影 $p\cos\alpha$，
则火炮受到反向的沿斜面冲量 $-p\cos\alpha$。
忽略炮弹相对火炮质量造成的质量修正，沿斜面动量定理：
\[
0-Mu=Mg\sin\alpha\,\tau-p\cos\alpha .
\]
所以
\[
\boxed{\tau=
\frac{p\cos\alpha-M\sqrt{2g\ell\sin\alpha}}
 {Mg\sin\alpha}}.
\]
\excheck 物理上要求 $\tau\ge0$，即
$p\cos\alpha\ge M\sqrt{2g\ell\sin\alpha}$；
若不满足，不可能以给定脉冲停止炮身。
结果在理想脉冲模型下成立，未考虑实际炮身运动时滑轮或轮轨阻力。

\exproblem{162}{27--28}
质量为 $m$ 的子弹水平射入悬挂于两根等长细绳上的
质量 $M$ 的物体中并嵌入（图1.42），随后
悬挂系统向上摆至细绳偏离竖直 $\theta$。
设 $m\ll M$、绳长为 $\ell$。
求：（a）碰前子弹速率；
（b）子弹初始动能中转化为热的比例。
\exsolution
碰撞时间很短，沿水平方向动量守恒。
碰撞后子弹与物体组成整体，瞬时速度为 $V$：
\[
mv=(M+m)V.
\]
碰撞结束到摆至最高点，机械能守恒：
\[
\frac12(M+m)V^2=(M+m)g\ell(1-\cos\theta),
\]
故
\[
V=\sqrt{2g\ell(1-\cos\theta)}
=2\sqrt{g\ell}\sin\frac\theta2.
\]
因此
\[
\boxed{v=\frac{2(M+m)}m\sqrt{g\ell}
\sin\frac\theta2
\simeq\frac{2M}{m}\sqrt{g\ell}\sin\frac\theta2}.
\]
碰撞中的机械能损失比例为
\[
\eta_{\rm heat}=
1-\frac{\tfrac12(M+m)V^2}{\tfrac12mv^2}
=1-\frac{m}{M+m}
=\boxed{\frac{M}{M+m}\simeq1-\frac mM}.
\]
\excheck 短时碰撞\emph{只守恒水平动量}，不守恒动能；
碰后摆起阶段才使用机械能守恒。

\exproblem{163}{28}
质量为 $M$ 的可水平滑动大物块置于光滑水平面，
其上放一质量为 $m$ 的小圆片，
圆片获得初始水平速率 $v$，并沿大物块的
光滑弧形表面上爬（图1.43）。
求圆片能够达到的最大高度（相对初始高度）。
\exsolution
水平外力为零，因此系统水平动量
$P_x=mv$ 守恒。初始总动能为
$K_0=\tfrac12mv^2$。
在圆片达到最大可能高度时，系统动能不能小于
同一水平动量对应的质心整体平移动能
$K_{\rm CM}=P_x^2/[2(M+m)]$。
若在上升最高位置的相对速度瞬时为零，
该下界可以达到，因而
\[
mgh_{\max}=
\frac12mv^2-\frac{m^2v^2}{2(M+m)}
=\frac{mMv^2}{2(M+m)}.
\]
于是
\[
\boxed{h_{\max}=\frac{Mv^2}{2g(M+m)}}.
\]
\excheck 这是原书理想光滑导轨装置可达到的极限高度；
$M\to\infty$ 时回到固定斜面上的 $v^2/(2g)$，
$m\gg M$ 时可达到的高度趋小。

\exproblem{164}{28}
质量为 $m$ 的小圆片从高度 $h$ 的光滑小山
由静止滑下，到达底部后进入质量为 $M$、
初始静止的水平木板（地面光滑），
圆片与木板之间存在摩擦，最终两者共同运动（图1.44）。
求：（1）此过程中摩擦力所做总功；
（2）这一结果是否依赖惯性参考系的选择。
\exsolution
圆片到山脚时速率
$u=\sqrt{2gh}$，木板静止。
进入木板后系统水平方向没有外力，
因此最终共同速度由动量守恒给出
\[
V=\frac{mu}{m+M}.
\]
摩擦耗散的总机械能等于最终与初始动能之差：
\[
W_f=\frac12(M+m)V^2-\frac12mu^2
=-\frac12\frac{mM}{m+M}u^2.
\]
用 $u^2=2gh$，得到
\[
\boxed{W_f=-\frac{mM}{m+M}gh=-\mu gh},
\qquad \mu=\frac{mM}{m+M}.
\]
（2）两物体的摩擦力大小相等、方向相反，故它们做功的
总和可以写为
$W_f=-\int f\,\mathrm d(x_{\rm disc}-x_{\rm plank})$。
两物体的相对位移在所有伽利略惯性系中相同，
所以 $\boxed{W_f\text{ 不依赖惯性系选择}}$。
\excheck 这里不变的是\emph{摩擦力的合计总功}，
并非摩擦力分别对两物体所做的单独功。

\exproblem{165}{28}
一颗石子从地面上方高度 $h$ 的地方相对地面静止释放，
忽略空气阻力。试改用一个相对地面以恒速
$v_0$ 向下运动的惯性参考系，重新得到
石子撞地速率 $v_0=\sqrt{2gh}$。
\exsolution
设所选参考系向下匀速平移，速率等于石子
\emph{最终}撞地时相对地面的速率 $v_0$。
在该系中，石子起始速度竖直向上，大小为 $v_0$；
由于重力向下，石子作加速度 $-g$ 的运动，
恰在撞地时相对该系速度为零。
这段时间 $T=v_0/g$，石子在移动系中向上位移
\[
\Delta y' =v_0T-\frac12gT^2=\frac{v_0^2}{2g}.
\]
地面在该移动系中向上走过距离
$v_0T=v_0^2/g$。
初始石子与地面相隔 $h$，撞地时位置重合，
故地面比石子多向上走过的距离恰为 $h$：
\[
h=\frac{v_0^2}{g}-\frac{v_0^2}{2g}
=\frac{v_0^2}{2g}
\quad\Longrightarrow\quad
\boxed{v_0=\sqrt{2gh}}.
\]
\excheck 这个参考系虽向下运动，却仍是\emph{惯性系}；
不能把地面在该系中的运动误当成零。

\exproblem{166}{28}
一质量为 $m_1=1.0\,\mathrm g$ 的质点
初速度 $\bm v_1=3.0\bm i-2.0\bm j\,(\mathrm{m/s})$，
与质量为 $m_2=2.0\,\mathrm g$、
初速度 $\bm v_2=4.0\bm j-6.0\bm k\,(\mathrm{m/s})$
的质点发生完全非弹性碰撞并结合。
求结合体的速度矢量及速率。
\exsolution
忽略外界冲量，总动量守恒：
\[
(m_1+m_2)\bm v=m_1\bm v_1+m_2\bm v_2.
\]
按各坐标分量计算
\[
\bm v=
\frac{1(3,-2,0)+2(0,4,-6)}{1+2}
=(1,2,-4)\,\mathrm{m/s}.
\]
因此
\[
\boxed{\bm v=(1.0\bm i+2.0\bm j-4.0\bm k)\,\mathrm{m/s}},
\qquad
\boxed{|\bm v|=\sqrt{1+4+16}=\sqrt{21}
\simeq4.58\,\mathrm{m/s}}.
\]
\excheck 分量计算使用相同质量单位 g，
故分子和分母的单位约去；碰撞后两球必有同一速度。

\exproblem{167}{28}
由质量 $m_1,m_2$ 的两球构成的孤立系统
发生完全非弹性碰撞，碰前速度分别为
$\bm v_1,\bm v_2$。求系统总动能增量。
\exsolution
碰后结合体速度由动量守恒确定：
\[
\bm V=\frac{m_1\bm v_1+m_2\bm v_2}{m_1+m_2}.
\]
动能变化为
\[
\Delta K=\frac12(m_1+m_2)V^2-
\frac12m_1v_1^2-\frac12m_2v_2^2.
\]
展开平方并合并同类项，得到
\[
\boxed{\Delta K=
-\frac12\mu|\bm v_1-\bm v_2|^2},
\qquad
\mu=\frac{m_1m_2}{m_1+m_2}.
\]
\excheck 动能变化非正，仅当两球碰前速度矢量完全相同
时为零。动量守恒不意味着动能守恒。

\exproblem{168}{28}
质量 $m_1$ 的粒子与静止质量 $m_2$ 的粒子
发生完全弹性碰撞。求入射粒子损失的初始动能比例：
（a）碰后入射粒子速度与原方向垂直；
（b）正碰。
\exsolution
设入射初速为 $v$，碰后入射速率为 $u$，
目标粒子碰后速度由动量守恒
$\bm w=(m_1/m_2)(\bm v-\bm u)$ 确定。
（a）若 $\bm v\cdot\bm u=0$，则
$w^2=(m_1/m_2)^2(v^2+u^2)$。
代入弹性碰撞动能守恒并整理：
\[
m_1v^2=m_1u^2+
\frac{m_1^2}{m_2}(v^2+u^2)
\Longrightarrow
\frac{u^2}{v^2}=\frac{m_2-m_1}{m_1+m_2}.
\]
故
\[
\boxed{\eta_a=1-\frac{u^2}{v^2}
=\frac{2m_1}{m_1+m_2}}.
\]
（b）一维正碰后入射粒子速度
$u=(m_1-m_2)v/(m_1+m_2)$，所以
\[
\boxed{\eta_b=
1-\left(\frac{m_1-m_2}{m_1+m_2}\right)^2
=\frac{4m_1m_2}{(m_1+m_2)^2}}.
\]
\excheck （a）要求 $m_2>m_1$ 才有非零速率的严格
$90^\circ$ 散射；$m_2=m_1$ 的极限入射粒子停下，
其“速度方向”不再有定义。不能让能量损失比例超过 $1$。

\exproblem{169}{28--29}
粒子1与静止粒子2发生完全弹性碰撞。求质量比 $m_1/m_2$：
（a）正碰后两粒子反向飞开且速率相等；
（b）碰后两粒子的飞行方向相对入射方向对称，
两方向夹角为 $\theta=60^\circ$。
\exsolution
（a）设碰前粒子1速度 $v$，
碰后两者速度为 $-u,+u$，其中 $u>0$。
动量与动能守恒分别给出
\[
m_1v=(m_2-m_1)u,\qquad
m_1v^2=(m_1+m_2)u^2.
\]
消去 $u,v$，得到
$(m_2-m_1)^2=m_1(m_1+m_2)$，
物理允许解
$\boxed{m_1/m_2=1/3}$。
（b）两射出方向分别相对初始方向倾斜
$+\theta/2$ 与 $-\theta/2$。
由横向动量守恒，
$m_1u_1=m_2u_2$；
由纵向动量守恒，
$m_1v=(m_1u_1+m_2u_2)\cos(\theta/2)
=2m_1u_1\cos(\theta/2)$。
因此 $u_1=v/[2\cos(\theta/2)]$。
动能守恒
$m_1v^2=m_1u_1^2+m_2u_2^2
=m_1u_1^2(1+m_1/m_2)$，
故
\[
\boxed{\frac{m_1}{m_2}=4\cos^2(\theta/2)-1
=1+2\cos\theta=2.0}.
\]
\excheck （b）要求两侧角度对称而非两粒子碰后速率相等；
横向动量守恒决定了两速率的质量加权关系。

\exproblem{170}{29}
质量相等的两光滑小球发生弹性碰撞，
其中一个小球以速率 $v$ 平动，另一个原先静止。
碰撞瞬间球心连线与入射速度成
$\alpha=45^\circ$ 角。
求两球压缩达到最大形变瞬间，
初始动能中暂时转为弹性势能的比例。
\exsolution
把入射速度分解为沿碰撞法线（球心连线）
的分量 $v_n=v\cos\alpha$，
及切向分量 $v_\tau=v\sin\alpha$。
光滑球无切向冲量，因此只在法向发生压缩。
在最大压缩瞬间，两球法向速度相等；
由法向动量守恒，两球共同法向速度为 $v_n/2$。
初始法向动能为 $\tfrac12mv_n^2$，
最大压缩时法向平动动能为
$2(\tfrac12 m(v_n/2)^2)=\tfrac14mv_n^2$，
故暂存弹性势能为
\[
U_{\rm elastic,max}=\frac14mv_n^2
=\frac14mv^2\cos^2\alpha .
\]
与初始总动能 $\tfrac12 mv^2$ 相比：
\[
\boxed{\eta=
\frac{U_{\rm elastic,max}}{\tfrac12 mv^2}
=\frac12\cos^2\alpha
=\frac14=25\%}.
\]
\excheck 完全弹性碰撞后的最终势能又恢复为零；
这里所问的是\emph{最大压缩瞬间}，不是碰后耗散热的比例。
% === END 03-mechanics-conservation-161-170.tex ===

% === BEGIN 03-mechanics-conservation-171-180.tex blob 967d489476b44c7221c8cefc22b62b016f26a4b0 ===
% Irodov (user-provided Arihant PDF), printed pp.29--30, PDF pp.38--39.
% Original problem numbers are preserved. Calculations independently expanded.
\exproblem{171}{29}
一炮弹以 $v=500\,\mathrm{m/s}$ 飞行，突然炸成质量相等的三块。
爆炸后三块的总动能是爆炸前的 $\eta=1.5$ 倍。
忽略爆炸瞬间外力冲量，问其中一块的最大可能速率。
\exsolution
设原总质量为 $M$，其中一块速度为 $\bm u$，
另两块速度为 $\bm v_2,\bm v_3$，爆炸前速度为 $\bm v$。
动量守恒给出
$\bm u+\bm v_2+\bm v_3=3\bm v$。
对于固定的 $\bm u$，另外两块的动能和在
$\bm v_2=\bm v_3=(3\bm v-\bm u)/2$ 时最小，
因为平方和不小于两倍平均矢量模平方。于是
\[
 K_f\ge\frac M6\left[
 u^2+\frac12|3\bm v-\bm u|^2\right]
 =\frac M2v^2+\frac M4|\bm u-\bm v|^2 .
\]
结合 $K_f=\eta Mv^2/2$ 得
$|\bm u-\bm v|^2\le2(\eta-1)v^2$。
三角不等式及其可取等条件给出
\[
\boxed{u_{\max}=v\bigl[1+\sqrt{2(\eta-1)}\bigr]
=1000\,\mathrm{m/s}}.
\]
最大值在该碎片与原速度同向、另两碎片同速时取得。
\excheck 对 $\eta=1$，没有额外爆炸动能时上界回到 $v$；
速度是矢量，单用一维能量守恒而不考虑动量约束会高估结果。

\exproblem{172}{29}
质量相等的两质点发生一维正碰。质点1碰前速率
$v=10\,\mathrm{m/s}$，质点2静止；碰后总动能损失
初始动能的 $\eta=1.0\%$。求碰后质点1的速度（含方向）。
\exsolution
取入射方向为正，碰后速度为 $u_1,u_2$。
动量守恒给出 $u_1+u_2=v$；碰后总动能为
$(1-\eta)mv^2/2$，因此
\[
u_1^2+(v-u_1)^2=(1-\eta)v^2,
\]
\[
u_1=\frac v2\left[1\pm\sqrt{1-2\eta}\right].
\]
碰后两球必须分离，因而 $u_2\ge u_1$；
选用减号分支，得到
\[
\boxed{u_1=\frac v2(1-\sqrt{1-2\eta})
\simeq0.0503\,\mathrm{m/s}},
\]
方向仍与原入射方向相同；$u_2\simeq9.9497\,\mathrm{m/s}$。
\excheck 加号分支意味着后一球比前一球慢，碰后仍会相互接近，
不代表一次完成的分离碰撞；题设要求 $0\le\eta\le1/2$。

\exproblem{173}{29}
质量 $m$ 的入射粒子碰前速率为 $v$，与静止质量
$M$ 的粒子碰撞后偏转 $90^\circ$；靶粒子的反冲方向与原
入射方向成 $\theta=30^\circ$，且 $M/m=5.0$。
求系统总动能的百分变化。
\exsolution
取原入射速度沿 $x$ 轴，入射粒子末速度沿 $+y$，
靶粒子反冲必在 $-y$ 一侧。设末速分别为 $u,V$，
两个方向的动量守恒为
\[
mv=MV\cos\theta,\qquad mu=MV\sin\theta.
\]
因此 $V=mv/(M\cos\theta)$、
$u=v\tan\theta$。末初动能比
\[
\frac{K_f}{K_i}
=\frac{mu^2+MV^2}{mv^2}
=\tan^2\theta+\frac mM\sec^2\theta.
\]
所以
\[
\boxed{\frac{\Delta K}{K_i}
=\left(1+\frac mM\right)\tan^2\theta+
\frac mM-1=-0.40}.
\]
\excheck 系统总动能减少 $40\%$，不是增加；此推导只要求
碰撞期间总动量守恒，不需要假定弹性碰撞。

\exproblem{174}{29}
孤立系统内两个质点质量为 $m_1,m_2$，
原参考系中的速度 $\bm v_1,\bm v_2$ 互相垂直。
求两质点在质心参考系中的动量及该系总动能。
\exsolution
质心速度
$\bm V=(m_1\bm v_1+m_2\bm v_2)/(m_1+m_2)$。
定义约化质量 $\mu=m_1m_2/(m_1+m_2)$，则
\[
\widetilde{\bm p}_1
=m_1(\bm v_1-\bm V)=\mu(\bm v_1-\bm v_2),
\quad
\widetilde{\bm p}_2=-\widetilde{\bm p}_1.
\]
由 $\bm v_1\cdot\bm v_2=0$，
\[
\boxed{|\widetilde{\bm p}_1|
=|\widetilde{\bm p}_2|=\mu\sqrt{v_1^2+v_2^2}},
\qquad
\boxed{\widetilde K
=\frac{\mu}{2}|\bm v_1-\bm v_2|^2
=\frac{\mu}{2}(v_1^2+v_2^2)}.
\]
\excheck 质心系总动量为零，但总动能通常不为零；
这里没有把两个粒子的\emph{原}动量当作质心系动量。

\exproblem{175}{29}
质量 $m_1>m_2$ 的粒子撞击静止质量 $m_2$ 的粒子，
碰撞完全弹性。求入射粒子在碰后相对原入射方向的
最大偏转角。
\exsolution
设入射速率 $v$，碰后入射粒子速度矢量 $\bm u$。
消去靶粒子的末速度
$\bm w=(m_1/m_2)(\bm v-\bm u)$，
由动能守恒得
\[
(m_1+m_2)u^2-2m_1vu\cos\theta+(m_1-m_2)v^2=0.
\]
把它视为关于 $u$ 的二次方程，要求存在实正根，
判别式非负：
\[
m_1^2\cos^2\theta-(m_1+m_2)(m_1-m_2)\ge0,
\quad \sin^2\theta\le\frac{m_2^2}{m_1^2}.
\]
所以
\[
\boxed{\theta_{\max}=\arcsin\frac{m_2}{m_1}}.
\]
\excheck $m_1\gg m_2$ 时偏转角很小；
当 $m_1\downarrow m_2$ 时最大偏转趋于 $90^\circ$。

\exproblem{176}{29}
三块质量相同的光滑圆盘 $A,B,C$ 在光滑水平面上。
$B,C$ 静止、两盘心间距等于盘直径的 $\eta$ 倍；
$A$ 沿两盘心连线的垂直平分线前进，速率 $v$，
同时与 $B,C$ 完全弹性碰撞（原图1.45）。
求 $A$ 碰后速度与反弹、停止、继续前进的判据。
\exsolution
盘半径设为 $R$，同时接触要求 $1\le\eta<2$。
碰撞时 $\overline{AB},\overline{AC}$ 与入射方向
的夹角为 $\phi$，几何给出
$\sin\phi=(2R\eta/2)/(2R)=\eta/2$，所以
$\cos^2\phi=1-\eta^2/4$。
设对 $B,C$ 的法向冲量大小均为 $J$，对称性使
$A$ 碰后仍仅沿原直线，速度有符号值 $u$：
\[
m(v-u)=2J\cos\phi.
\]
两目标盘碰后速度大小都为 $J/m$；
由于碰撞弹性，
\[
mv^2=mu^2+2J^2/m.
\]
解以上两式、舍去无碰撞的零冲量分支，得到
\[
\boxed{u=v\frac{\eta^2-2}{6-\eta^2}}.
\]
因此 $1\le\eta<\sqrt2$ 时 $A$ 反弹；
$\eta=\sqrt2$ 时停止；
$\sqrt2<\eta<2$ 时继续向前。
\excheck 这里的 $u$ 是\emph{有符号速度}；图形的双接触
只在指定中心距区间有意义，不能延伸到 $\eta>2$。

\exproblem{177}{29--30}
一个分子与另一个同质量、碰前静止的分子碰撞。
证明：（a）碰撞理想弹性时，两粒子的末速度方向夹角
为 $90^\circ$；（b）若动能损失为正，夹角不再是 $90^\circ$。
\exsolution
两粒子质量均为 $m$，碰前入射速度为 $\bm v$，
末速度为 $\bm u,\bm w$。动量守恒
$\bm v=\bm u+\bm w$，故
\[
v^2=u^2+w^2+2\bm u\cdot\bm w.
\]
若完全弹性，动能守恒
$v^2=u^2+w^2$，所以
$2\bm u\cdot\bm w=0$；当二者末速都非零时，
\[
\boxed{\angle(\bm u,\bm w)=90^\circ}.
\]
若发生非弹性碰撞且动能损失 $\Delta E>0$，
则 $v^2>u^2+w^2$，故
\[
2\bm u\cdot\bm w
=v^2-u^2-w^2=2\Delta E/m>0,
\]
夹角为锐角。
\excheck 如有一粒子恰停下，其速度方向未定义；
“夹角为 $90^\circ$”指两末速度非零的通常散射情形。

\exproblem{178}{30}
火箭在某时刻质量为 $m$、加速度为 $\bm a$，
外力合力为 $\bm F$。火箭以恒定质量流率
$\mu>0$（$\mathrm{kg/s}$）喷出气体，
气体相对火箭的速度矢量为 $\bm u$。
证明火箭方程 $m\bm a=\bm F-\mu\bm u$。
\exsolution
取时间间隔 $\mathrm dt$，火箭初始速度 $\bm v$，
其质量减少 $\mu\,\mathrm dt$；
喷出的气体质量为 $\mu\,\mathrm dt$，
在惯性系中速度为 $\bm v+\bm u$。
动量变化保留到 $\mathrm dt$ 一阶：
\[
\begin{aligned}
\Delta\bm P
&=(m-\mu\mathrm dt)(\bm v+\mathrm d\bm v)
 +\mu\mathrm dt(\bm v+\bm u)-m\bm v\\
&=m\,\mathrm d\bm v+\mu\bm u\,\mathrm dt
+O(\mathrm dt^2).
\end{aligned}
\]
由外力冲量 $\bm F\,\mathrm dt=\Delta\bm P$，
\[
\boxed{m\bm a=\bm F-\mu\bm u}.
\]
\excheck 此处 $\bm u$ 是\emph{喷气相对火箭的矢量速度}，
通常指向火箭后方；反冲推力方向为 $-\mu\bm u$。

\exproblem{179}{30}
火箭初始静止、质量 $m_0$，在无外力情况下
以相对火箭不变的喷气速度 $\bm u$ 排气。
求火箭质量变为 $m$ 时的速度。
\exsolution
令 $\mu=-\mathrm dm/\mathrm dt>0$，
用上一题方程且 $\bm F=\bm0$：
\[
m\frac{\mathrm d\bm v}{\mathrm dt}
=-\mu\bm u=\bm u\,\frac{\mathrm dm}{\mathrm dt}.
\]
两端分离变量并由初态积分：
\[
\int_{\bm0}^{\bm v}\mathrm d\bm v
=\bm u\int_{m_0}^{m}\frac{\mathrm dm'}{m'},
\]
所以
\[
\boxed{\bm v=\bm u\ln\frac m{m_0}
=-\bm u\ln\frac{m_0}{m}}.
\]
\excheck $m<m_0$，若 $\bm u$ 指向火箭后方，
速度 $\bm v$ 指向前方；其大小是
$|\bm u|\ln(m_0/m)$，与喷气质量流率无关。

\exproblem{180}{30}
在无外力作用下，火箭保持大小为 $a>0$ 的恒加速度，
喷气相对火箭的反向速率固定为 $u>0$，
初质量为 $m_0$。求质量 $m(t)$。
\exsolution
设推进方向为正，气体相对火箭速度为 $-u$。
火箭方程 $ma=-\mu(-u)=\mu u$，
其中 $\mu=-\dot m$，故
\[
ma=-u\frac{\mathrm dm}{\mathrm dt},
\qquad\frac{\mathrm dm}{m}=-\frac au\,\mathrm dt.
\]
由 $m(0)=m_0$ 得
\[
\boxed{m(t)=m_0e^{-at/u}}.
\]
\excheck 指数 $at/u$ 无量纲；质量单调减少且始终为正。
恒加速度要求质量流率按 $\mu(t)=am(t)/u$
随时间衰减，并非恒定流率。
% === END 03-mechanics-conservation-171-180.tex ===

% === BEGIN 03-mechanics-conservation-181-190.tex blob 2d2d8b0d7584e3958757c123f26f665f29005264 ===
% Irodov user-provided PDF, original problems 181--190, printed pp.30--32.
\exproblem{181}{30}
一飞船在无外力下最初以恒速 $v_0$ 飞行。为改变方向，
开启喷气发动机，始终使喷气相对飞船的速度大小为 $u$、
方向垂直于\emph{当时}飞船速度。飞船质量由 $m_0$ 减到 $m$
即停止喷气。求飞行方向的总偏转角。
\exsolution
喷气推力大小 $u\mu$，与飞船瞬时速度垂直，
故 $m\dot v=0$，飞船速率始终为 $v_0$。
推力在速度法线方向产生
\[
m v_0\frac{\mathrm d\alpha}{\mathrm dt}=u\mu,
\qquad \mu=-\frac{\mathrm dm}{\mathrm dt}>0.
\]
于是 $\mathrm d\alpha=-(u/v_0)\,\mathrm dm/m$，
从质量 $m_0$ 积分到 $m$，得到
\[
\boxed{\alpha=\frac{u}{v_0}\ln\frac{m_0}{m}}.
\]
\excheck 角度采用弧度；此结果要求喷气方向始终跟随飞船
当前运动方向转动，不是固定与最初速度垂直。

\exproblem{182}{30}
一辆装沙的小车在光滑水平面上，受与前进方向相同的恒力
$F$。沙从车底孔以恒定质量流率 $\mu>0$ 漏出；
漏出时相对车没有水平速度。最初小车与沙总质量
$m_0$、初速零。求任意 $0\le t<m_0/\mu$ 时的加速度、速度。
\exsolution
漏出的沙在离车一刻具有与车相同的水平速度，
因此没有相对水平喷流推力。车与剩余沙总质量
$m(t)=m_0-\mu t$，水平运动方程
\[
m(t)\frac{\mathrm dv}{\mathrm dt}=F.
\]
直接可得
\[
\boxed{a(t)=\frac{F}{m_0-\mu t}},\qquad
\boxed{v(t)=\int_0^t\frac{F\,\mathrm ds}{m_0-\mu s}
=\frac F\mu\ln\frac{m_0}{m_0-\mu t}}.
\]
\excheck 结果仅在沙尚未漏完的阶段成立；当剩余质量趋零的
数学发散不代表真实小车质量会趋零，应在实际空车质量处终止模型。

\exproblem{183}{30}
一质量为 $m_0$ 的空平板车最初静止，
受沿水平方向的恒力 $F$。静止料斗以恒质量流率
$\mu>0$ 向车中注入沙子（图1.46），沙子落入前
水平速度为零，摩擦忽略。求装料阶段速度与加速度的时间规律。
\exsolution
在固定惯性系中，车与已装入沙的质量为
$m(t)=m_0+\mu t$。
进入系统的沙子带入的水平动量通量为零，
故变质量系统的水平动量平衡为
\[
\frac{\mathrm d}{\mathrm dt}[m(t)v(t)]=F.
\]
由初始总水平动量零，积分得到
$m(t)v(t)=Ft$，于是
\[
\boxed{v(t)=\frac{Ft}{m_0+\mu t}},\qquad
\boxed{a(t)=\frac{\mathrm dv}{\mathrm dt}
=\frac{Fm_0}{(m_0+\mu t)^2}}.
\]
\excheck 此题与上一题不同：正在装入的沙子在地面系
水平静止，必须由车加速，不能错误地套用 $m\dot v=F$。

\exproblem{184}{31}
图1.47中，一根长度为 $\ell$ 的均匀柔软链条置于光滑
水平管内，一段竖直下垂长度为 $h$，其末端恰触桌面。
释放管中另一端 $A$，链条开始从管中滑出，
下垂的链条逐渐在桌面堆积。求 $A$ 刚滑出管口的速率。
\exsolution
设 $x$ 是从管内滑出的累计长度。
因竖直管口至桌面高度固定为 $h$，
已经堆积在桌面上的链长为 $x$，
仍在运动的链长为 $\ell-x$。
设线密度为 $\lambda$，运动链条的重力驱动力
始终为 $\lambda gh$。
在桌面上终止运动的链元携走与当前链条相同的动量，
对运动链段应用含质量流出的动量平衡，通量项抵消
质量随时间变化的相应项，得到
\[
\lambda(\ell-x)\frac{\mathrm dv}{\mathrm dt}
=\lambda gh .
\]
再用 $\dot x=v$，得
\[
v\frac{\mathrm dv}{\mathrm dx}=\frac{gh}{\ell-x}.
\]
从 $v(0)=0$ 积分到 $x=\ell-h$（此时 $A$ 恰到管口）：
\[
\frac{v^2}{2}
=gh\int_0^{\ell-h}\frac{\mathrm dx}{\ell-x}
=gh\ln\frac{\ell}{h}.
\]
所以
\[
\boxed{v=\sqrt{2gh\ln\frac{\ell}{h}}}.
\]
\excheck 不能直接把重力势能损失全部算成末态动能；
链元落桌后非弹性地停住，存在相应动能损失。
需有 $0<h<\ell$，当 $h\to\ell$ 时结果趋零。

\exproblem{185}{31}
某质点相对固定点 $O$ 的角动量矢量随时间变化为
$\bm L(t)=\bm a+\bm b\,t^2$，其中恒定矢量
$\bm a\perp\bm b$。当关于 $O$ 的力矩
$\bm N$ 与角动量夹角为 $45^\circ$ 时，
求力矩大小。
\exsolution
由角动量定理
\[
\bm N=\frac{\mathrm d\bm L}{\mathrm dt}=2t\,\bm b.
\]
设 $a=|\bm a|,\ b=|\bm b|$。
由于 $\bm a\perp\bm b$，
力矩方向沿 $\bm b$（取 $t>0$）；
$\bm L$ 沿 $\bm b$ 的分量大小为 $bt^2$，
垂直分量大小为 $a$。两矢量夹角为 $45^\circ$ 意味着
$bt^2=a$，故
\[
\boxed{|\bm N|=2bt=2\sqrt{ab}}.
\]
\excheck $[\bm b]=[\bm L]/\mathrm s^2$，
所以 $\sqrt{ab}$ 具有力矩量纲；
$t<0$ 时只需用 $|t|$，大小结论不变。

\exproblem{186}{31}
质量为 $m$ 的球以初速度 $v_0$、仰角 $\alpha$
从地面某点抛出，忽略空气阻力。求球对抛出点的
角动量大小随时间变化的规律；并求
$m=0.130\,\mathrm{kg}$、$\alpha=45^\circ$、
$v_0=25\,\mathrm{m/s}$ 时在轨迹顶点的角动量大小。
\exsolution
取抛出点为原点，$x$ 轴水平、$y$ 轴向上：
\[
\bm r=(v_0\cos\alpha\,t,\,
v_0\sin\alpha\,t-\tfrac12gt^2),\quad
\bm v=(v_0\cos\alpha,\,v_0\sin\alpha-gt).
\]
角动量只有垂直于运动平面的分量：
\[
L_z=m(xv_y-yv_x)=-\frac12mgv_0t^2\cos\alpha.
\]
因此
\[
\boxed{|\bm L(t)|=\frac12mgv_0t^2\cos\alpha}.
\]
顶点时间 $t_h=v_0\sin\alpha/g$，
\[
\boxed{L_h=\frac{mv_0^3}{2g}\sin^2\alpha\cos\alpha
\simeq36.6\,\mathrm{kg\,m^2/s}}.
\]
\excheck 取 $g=9.81\,\mathrm{m/s^2}$ 时约 $37$，
与底本答案的两位有效数字相符。角动量矢量方向按
所选坐标为负 $z$ 方向，题目只要求模。

\exproblem{187}{31}
光滑圆片以速率 $v$ 沿光滑水平面滑向固定光滑墙，
碰撞为理想弹性；入射方向与墙面法线夹角为 $\alpha$，
碰撞点为 $O$（图1.48）。
求：（a）相对哪些点碰撞前后角动量守恒；
（b）墙面内、在运动平面上距 $O$ 为 $\ell$ 的点 $O'$
处，圆片角动量变化的模。
\exsolution
碰撞中墙对圆片的冲量仅沿墙面法线。
弹性反射使法线速度分量反号，切线速度分量不变，
故动量增量大小为
\[
|\Delta\bm p|=2mv\cos\alpha,
\]
方向沿墙面法线。关于任意固定点 $Q$ 的角动量变化
是冲量矩
$\Delta\bm L_Q=(\bm r_O-\bm r_Q)\times\Delta\bm p$。
因此当 $Q$ 位于\emph{过 $O$ 的墙面法线直线}上，
力臂与冲量平行，角动量保持不变。
对于题给的 $O'$，其与 $O$ 之间的线段沿墙，
垂直于冲量，于是
\[
\boxed{|\Delta\bm L_{O'}|
=\ell|\Delta\bm p|=2mv\ell\cos\alpha}.
\]
\excheck 虽然关于撞击点 $O$ 的角动量不变，
关于墙上另一点 $O'$ 一般并不不变；
角动量守恒必须说明参考点。

\exproblem{188}{31}
小球质量 $m$，用长为 $\ell$ 的轻绳悬挂于固定点 $O$，
绕竖直轴以恒角速度 $\omega$ 作水平圆周运动（圆锥摆）。
问关于哪个点角动量矢量守恒；
再求关于绳悬点 $O$ 的角动量矢量在半周内的变化量大小。
\exsolution
记小球运动圆半径为 $\rho=\ell\sin\theta$，
圆平面低于 $O$ 的距离为 $h=\ell\cos\theta$。
关于圆心 $C$ 的位置与动量都在水平面内，
角动量始终为沿竖直轴的恒矢量
$\bm L_C=m\omega\rho^2\bm e_z$。
关于 $O$，将位置矢量写为
$\bm r_{O}=\rho\bm e_r-h\bm e_z$，速度为
$\bm v=\omega\rho\bm e_\varphi$，
所以
\[
\bm L_O
=m\omega\rho^2\bm e_z+
m\omega\rho h\,\bm e_r .
\]
一段半周后 $\bm e_r\mapsto-\bm e_r$，
竖直分量不变，故
\[
|\Delta\bm L_O|=2m\omega\rho h.
\]
圆锥摆的动力学关系为
$\omega^2\ell\cos\theta=g$，
最终
\[
\boxed{|\Delta\bm L_O|
=\frac{2mg\ell}{\omega}
\sqrt{1-\left(\frac g{\omega^2\ell}\right)^2}}.
\]
\excheck 只要 $\omega^2\ell\ge g$ 才有非退化圆锥摆；
题目所求“守恒的参考点”为\emph{运动圆的圆心} $C$。

\exproblem{189}{31}
质量 $m$ 的球从高度 $h$、地面系静止释放，
忽略空气阻力。考察以恒定水平速度
$\bm V$ 平移的惯性系，取球的释放点在该系初始为 $O$。
求自由落体期间关于 $O$ 的角动量矢量增量的模。
\exsolution
设地面系竖直向下为 $y$ 轴，
平移系水平方向为 $x$ 轴。
在平移系中球的速度与位置为
\[
\bm v'=(-V,\,gt),\qquad
\bm r'=(-Vt,\,\tfrac12gt^2).
\]
垂直于运动平面的角动量分量：
\[
L_z'=m(x'v_y'-y'v_x')
=m(-Vgt^2+\tfrac12Vgt^2)
=-\tfrac12mVgt^2.
\]
初始为零，落地时间 $t_f=\sqrt{2h/g}$，
从而
\[
\boxed{|\Delta\bm L|=\frac12mVg\,t_f^2=mVh}.
\]
\excheck 使用运动参考系中\emph{固定}于该系的点 $O$；
关于地面系中的释放点，球垂直下落时角动量始终为零，
两者并不矛盾。

\exproblem{190}{31--32}
水平光滑圆盘绕中心固定竖轴以角速度 $\omega$ 匀速转动，
$t=0$ 时，一小圆片从盘中心以相对盘初速
$v_0$ 向外滑行，盘面不施加水平力。
求在转盘参考系中关于盘心的角动量 $L_z(t)$，
并证明其变化由科里奥利力矩引起。
\exsolution
地面惯性系中圆片从原点沿直线以恒速 $v_0$ 运动，
所以 $r(t)=v_0t$。在转盘系中相对速度
$\bm v'=\bm v-\bm\omega\times\bm r$，
其切向分量为 $v_\varphi'=-\omega r$。
关于盘心的相对角动量为
\[
\boxed{L_z'(t)=mr v_\varphi'
=-m\omega v_0^2t^2},
\qquad |L_z'|=m|\omega|v_0^2t^2.
\]
旋转系中离心惯性力沿径向，关于盘心力矩为零；
科里奥利力 $\bm F_{\rm Cor}=-2m\bm\omega\times\bm v'$
沿竖直轴的力矩为
$N_z^{\rm Cor}=-2m\omega r\dot r
=-2m\omega v_0^2t$。
故
\[
\frac{\mathrm dL_z'}{\mathrm dt}
=-2m\omega v_0^2t=N_z^{\rm Cor}.
\]
\excheck 本版书末只写角动量\emph{模}，
而按 $\omega>0$（逆时针）约定的有符号结果为负；
二者一致。
% === END 03-mechanics-conservation-181-190.tex ===

% === BEGIN 03-mechanics-conservation-191-199.tex blob 7fe44718a6530d5ff8e03ba738d3644db562ca13 ===
% Irodov, source PDF pp.41--42; original printed pp.32--33.
\exproblem{191}{32}
质量 $m$ 的粒子在中心势场 $U(r)=kr^2$ 中作闭合轨道运动，
$k>0$。已知粒子与中心 $O$ 的最近距离为 $r_1$，
最远点处速率为 $v_2$，求 $m$。
\exsolution
在最近点和最远点，径向速度均为零，速率完全切向。
记最远半径为 $r_2$、守恒角动量为 $L$，则
$L=mr_1v_1=mr_2v_2$。
能量守恒可写成两端同一有效势：
\[
kr_1^2+\frac{L^2}{2mr_1^2}
=kr_2^2+\frac{L^2}{2mr_2^2}.
\]
对非圆形轨道 $r_2\ne r_1$，两边移项并约去
$r_2^2-r_1^2$，得到
$L^2/(2m)=kr_1^2r_2^2$。再由
$L=mr_2v_2$ 消去 $L,r_2$：
\[
\boxed{m=\frac{2kr_1^2}{v_2^2}}.
\]
\excheck 此结论来自中心力场的角动量与总能量同时守恒；
若轨道恰为圆，$r_2=r_1$ 时不能采用上述约因子步骤，
但圆轨道向心条件同样给出此式。

\exproblem{192}{32}
长为 $\ell$ 的轻绳悬挂一小球，
初始与竖直方向夹角为 $\theta$。在这一位置使小球
获得与其所在竖直平面垂直的水平初速 $v_0$。
求使小球运动中能达到最大偏角 $90^\circ$ 的临界初速。
\exsolution
绕竖直轴的外力矩为零，轴向角动量守恒。
初始绕轴半径为 $\ell\sin\theta$，所以
$L_z=m\ell v_0\sin\theta$。
当偏角达到 $\pi/2$ 的转折点时，
沿经线方向的速度瞬时为零，只有绕竖轴的方位速率 $u$，
此时半径为 $\ell$，故
$u=v_0\sin\theta$。
以初始点为势能零点，机械能守恒：
\[
\frac12mv_0^2
=\frac12m v_0^2\sin^2\theta+mg\ell\cos\theta.
\]
由此
\[
\boxed{v_{0,\mathrm{crit}}=
\sqrt{\frac{2g\ell}{\cos\theta}}}.
\]
\excheck 需 $0\le\theta<\pi/2$。
该临界速度对应“恰好到达 $90^\circ$”的转折状态，
初速更大将越过此角度。

\exproblem{193}{32}
质量为 $m$ 的小物体被不可伸长细绳系住，
在光滑水平面上绕固定孔 $O$ 转动。
有人以恒速将细绳向下拉入孔中，故物体到孔距离
$r$ 以恒定速率减小；在 $r=r_0$ 时角速度为 $\omega_0$。
求拉线张力 $T(r)$。
\exsolution
由于细绳张力始终径向指向孔，
关于孔 $O$ 的外力矩为零，所以
\[
mr^2\dot\varphi=mr_0^2\omega_0,\qquad
\dot\varphi=\frac{r_0^2\omega_0}{r^2}.
\]
在极坐标中径向加速度为
$a_r=\ddot r-r\dot\varphi^2$。
拉入速度恒定意味着 $\ddot r=0$，
径向牛顿第二定律 $-T=ma_r$ 得
\[
\boxed{T(r)=mr\dot\varphi^2
=\frac{m\omega_0^2r_0^4}{r^3}}.
\]
\excheck 张力随半径减小迅速增大。
若拉线速度并非常数，还必须保留 $-m\ddot r$ 项。

\exproblem{194}{32}
一根不可伸长轻绳绕在半径为 $R$ 的有质量定轴滑轮上，
自由端系一小物体质量为 $m$。从静止开始释放，
求\emph{滑轮和重物组成的整个系统}相对滑轮固定轴的
总角动量随时间的变化。
\exsolution
以滑轮固定转轴为矩心，支承轴反力对该轴不产生力矩；
绳张力是系统内部相互作用，其作用力矩在滑轮与重物的
总角动量平衡中抵消。
重物始终在距轴水平距离 $R$ 的铅垂线上运动，
其重力关于轴的力矩大小为 $mgR$，方向恒定。
对整个系统应用角动量定理：
\[
\frac{\mathrm dL_{\rm total}}{\mathrm dt}=mgR,
\qquad L_{\rm total}(0)=0,
\]
因此
\[
\boxed{|L_{\rm total}(t)|=mgRt}.
\]
\excheck 结果与滑轮转动惯量无关，因为求的是\emph{整个}
滑轮—重物系统的角动量。单独求重物的角动量不能在
未知滑轮转动惯量时直接写成 $mgRt$。

\exproblem{195}{32}
均匀实心球质量 $m$、半径 $R$，从静止开始沿
倾角为 $\alpha$ 的斜面无滑动滚下。
求相对于\emph{最初接触点} $O$ 的角动量大小 $L(t)$；
若斜面完全光滑，结果如何变化？
\exsolution
对于纯滚动均匀实心球，
$I_C=2mR^2/5$，质心沿斜面加速度
$a_C=(5/7)g\sin\alpha$，
速率为 $v=a_Ct$，角速度 $\omega=v/R$。
关于初始接触点分解：
\[
\bm L_O=(\bm r_C-\bm r_O)\times m\bm v_C
+I_C\bm\omega.
\]
球心相对初始接触点的沿斜面位移与 $\bm v_C$ 平行，
故其叉积贡献为零；垂直斜面的距离始终为 $R$，
于是两部分角动量同向，大小相加为
\[
L_O=Rm v+I_C\frac vR
=\frac75mRv=mgR\sin\alpha\,t.
\]
若斜面完全光滑，球只平动不自转，
$a_C=g\sin\alpha$，则
$L_O=Rm(g\sin\alpha\,t)$，因此
\[
\boxed{L_O(t)=mgR\sin\alpha\,t}
\]
在两种情况下相同。
\excheck 本题需要区分\emph{重力力矩}与
\emph{合外力矩}：选最初接触点$O$为固定参考点，
球心沿斜面已移动$s$，则重力相对$O$
的力臂并非一直等于$R\sin\alpha$。
取沿斜面下行为$\bm e_t$、离开斜面为$\bm e_n$，
球心位置$\bm r_C=s\bm e_t+R\bm e_n$。
重力和接触支持力（$N=mg\cos\alpha$）的
关于$O$力矩中，分别有正负
$s mg\cos\alpha$项，恰好抵消；
接触静摩擦力与$\bm r_{\rm contact}=s\bm e_t$
共线，不产生关于$O$的力矩。
因此\emph{总外力矩}的大小才恒为
$mgR\sin\alpha$，对有无摩擦均成立。
不能把“总外力矩恒定”误写为“重力力矩恒定”。

\exproblem{196}{32}
一个粒子系统相对于固定点 $O$ 的总动量为 $\bm p$，
总角动量为 $\bm L_O$。另一固定点 $O'$ 的位置矢量
相对 $O$ 为 $\bm r_0$。求系统相对 $O'$ 的角动量，
并问在何种情形下该量与参考点选取无关。
\exsolution
对第 $i$ 粒子，相对 $O'$ 的位置矢量为
$\bm r_i'=\bm r_i-\bm r_0$。因此
\[
\begin{aligned}
\bm L_{O'}
&=\sum_i(\bm r_i-\bm r_0)\times\bm p_i\\
&=\sum_i\bm r_i\times\bm p_i
-\bm r_0\times\sum_i\bm p_i.
\end{aligned}
\]
故
\[
\boxed{\bm L_{O'}=\bm L_O-\bm r_0\times\bm p}.
\]
任意点的改变都不影响角动量，当且仅当
$\boxed{\bm p=\bm0}$。
\excheck 对某一特定参考点改变，若 $\bm r_0\parallel\bm p$
也可能有相同角动量；但要求\emph{任意}选点均不变则必须 $\bm p=0$。

\exproblem{197}{32}
证明柯尼希角动量分解：系统相对于固定点 $O$ 的角动量为
$\bm L_O=\widetilde{\bm L}+\bm r_C\times\bm p$，
其中 $\widetilde{\bm L}$ 是质心平动参考系中的内禀角动量，
$\bm r_C$ 是质心相对 $O$ 的位置矢量，
$\bm p$ 为总动量。
\exsolution
写出每个质点的位置和速度相对质心的分解：
\[
\bm r_i=\bm r_C+\bm\rho_i,\qquad
\bm v_i=\bm V_C+\dot{\bm\rho}_i.
\]
质心定义保证
$\sum_i m_i\bm\rho_i=\bm0$ 和
$\sum_i m_i\dot{\bm\rho}_i=\bm0$。
展开角动量总和：
\[
\begin{aligned}
\bm L_O
&=\sum_i m_i(\bm r_C+\bm\rho_i)
\times(\bm V_C+\dot{\bm\rho}_i)\\
&=M\bm r_C\times\bm V_C+
\sum_i m_i\bm\rho_i\times\dot{\bm\rho}_i,
\end{aligned}
\]
因为两个交叉项均为零。
令 $\bm p=M\bm V_C$、
$\widetilde{\bm L}=
\sum_i m_i\bm\rho_i\times\dot{\bm\rho}_i$，得
\[
\boxed{\bm L_O=\widetilde{\bm L}+\bm r_C\times\bm p}.
\]
\excheck 此分解是恒等式，并不要求系统孤立，也不要求
内禀角动量本身守恒。

\exproblem{198}{33}
质量 $m$ 的小球以速率 $v_0$ 水平正碰一个初始静止的
刚性哑铃的一端（图1.50）。哑铃两端小球质量各为 $m/2$，
端点相距 $\ell$，忽略端球尺寸，
碰撞完全弹性。求碰后哑铃\emph{质心系}角动量大小。
\exsolution
哑铃总质量为 $m$，绕自身质心转动惯量
$I=(m/2)(\ell/2)^2+(m/2)(\ell/2)^2=m\ell^2/4$。
设水平碰撞冲量为 $J$；则哑铃质心平动速度
$V=J/m$，角速度满足
$I\omega=J\ell/2$，故
$\omega=2J/(m\ell)$。
受撞端点沿碰撞方向的速度为
$V+\omega\ell/2=2J/m$；
入射小球碰后速度为 $v_0-J/m$。
完全弹性意味着两接触点分离的相对速度等于
碰前接近速率：
\[
\frac{2J}{m}-\left(v_0-\frac Jm\right)=v_0,
\]
所以 $J=(2/3)mv_0$。
哑铃质心系角动量是该冲量对质心的矩：
\[
\boxed{\widetilde L=I\omega=\frac{J\ell}{2}
=\frac13m\ell v_0}.
\]
\excheck 这里哑铃的质心平动动能和转动动能都参与了
碰撞能量分配；不能把哑铃当作固定转轴的转子。

\exproblem{199}{33}
两块相同质量 $m$ 的小圆盘位于光滑水平面上，
用初始自然长 $\ell_0$、劲度系数 $\kappa$ 的轻弹簧连接。
使其中一块沿垂直于弹簧的方向获得速率 $v_0$。
若最大相对伸长 $\varepsilon_{\max}=
(\ell_{\max}-\ell_0)/\ell_0$ 很小，
求其近似值。
\exsolution
在质心系用相对坐标分析，约化质量
$\mu=m/2$。初始相对速度大小为 $v_0$，且与
连线垂直，所以相对角动量为
$L=\mu\ell_0v_0$。
初始能量为 $\mu v_0^2/2$。
当弹簧伸长到最大时径向速度为零，设
$r=\ell_0(1+\varepsilon)$，能量守恒：
\[
\frac12\mu v_0^2
=\frac{L^2}{2\mu r^2}
+\frac12\kappa\ell_0^2\varepsilon^2.
\]
代入 $L$ 并移项：
\[
\mu v_0^2\left[1-\frac1{(1+\varepsilon)^2}\right]
=\kappa\ell_0^2\varepsilon^2.
\]
若 $\varepsilon\ll1$，
$1-(1+\varepsilon)^{-2}=2\varepsilon+O(\varepsilon^2)$，
解得非零的首阶结果
\[
\boxed{\varepsilon_{\max}\simeq
\frac{m v_0^2}{\kappa\ell_0^2}}.
\]
\excheck 右端无量纲，且须自身远小于 $1$
才符合题目的小伸长近似条件。
% === END 03-mechanics-conservation-191-199.tex ===
